%% ============================================================================
%% Gợi ý tin tức sử dụng các kỹ thuật LLM (LLMRec) -- bản trình bày lại cho dễ theo dõi.
%% Biên dịch: pdfLaTeX, chạy 2 lần.
%%  * Nếu cùng thư mục có theme "hust" (beamerthemehust.sty) thì dùng theme đó;
%%    nếu không, dùng giao diện dự phòng định nghĩa ngay bên dưới.
%%  * Bản có lời thoại gợi ý (mỗi trang: slide ở trên, lời thoại ở dưới):
%%    bỏ dấu % ở hai dòng "pgfpages" và "show notes" trong phần khai báo.
%% Số liệu lấy từ các lần chạy đo thật: v7, v13 (MIND-small); v14, v15, v16b, v17 (MIND-large).
%% Các dòng đánh dấu (*) lấy từ lần chạy v17; các dòng còn lại của MIND-large lấy từ v16b.
%% Các dòng ghi "v20c" là mô hình mới nhất (v19c + thành phần LLM-text của v20a, v20b), đo ghép cặp trên cùng impression.
%% ============================================================================
\documentclass[aspectratio=169,t]{beamer}
\usepackage[utf8]{inputenc}
\usepackage[vietnamese]{babel}
\usepackage{amsmath,amssymb}
\usepackage{booktabs}
\usepackage{array}
\usepackage{tabularx}
\usepackage{multirow}
\usepackage{makecell}
\usepackage{graphicx}
\usepackage{tikz}
\usetikzlibrary{positioning,calc,decorations.pathreplacing,arrows.meta}
% \usepackage{pgfpages}
% \setbeameroption{show notes on second screen=bottom}

%% ---------------------------------------------------------------- màu dùng chung
\definecolor{okc}{HTML}{1E7B34}      % thành công
\definecolor{partc}{HTML}{B7791F}    % có điều kiện
\definecolor{failc}{HTML}{B91C1C}    % không hiệu quả
\definecolor{nac}{HTML}{64748B}      % chưa đo / chú thích
\definecolor{keyc}{HTML}{14284B}     % xanh đậm
\definecolor{softc}{HTML}{F2F4F8}
\definecolor{accentc}{HTML}{A6192E}  % đỏ Bách khoa

%% ---------------------------------------------------------------- theme
\newif\ifhusttheme
\IfFileExists{beamerthemehust.sty}{\hustthemetrue}{\hustthemefalse}
\ifhusttheme
  \usetheme{hust}
\else
  %% Giao diện dự phòng (chỉ dùng khi không có beamerthemehust.sty)
  \usetheme{default}
  \setbeamercolor{structure}{fg=accentc}
  \setbeamercolor{normal text}{fg=black!88}
  \setbeamercolor{frametitle}{fg=white,bg=accentc}
  \setbeamercolor{title}{fg=accentc}
  \setbeamercolor{subtitle}{fg=keyc}
  \setbeamercolor{block title}{fg=white,bg=keyc}
  \setbeamercolor{block body}{fg=black,bg=keyc!5}
  \setbeamercolor{block title alerted}{fg=white,bg=accentc}
  \setbeamercolor{block body alerted}{fg=black,bg=accentc!5}
  \setbeamercolor{block title example}{fg=white,bg=okc}
  \setbeamercolor{block body example}{fg=black,bg=okc!5}
  \setbeamercolor{foot}{fg=black!60,bg=black!5}
  \setbeamertemplate{blocks}[rounded][shadow=false]
  \setbeamertemplate{itemize items}[square]
  \setbeamertemplate{enumerate items}[default]
  \setbeamerfont{frametitle}{size=\large,series=\bfseries}
  \setbeamerfont{title}{size=\LARGE,series=\bfseries}
  \setbeamertemplate{frametitle}{%
    \nointerlineskip
    \begin{beamercolorbox}[wd=\paperwidth,ht=0.85cm,dp=0.3cm,leftskip=0.35cm,rightskip=0.35cm]{frametitle}%
      \usebeamerfont{frametitle}\insertframetitle
    \end{beamercolorbox}\vspace{1mm}}
  \setbeamertemplate{footline}{%
    \leavevmode\hbox{%
    \begin{beamercolorbox}[wd=.34\paperwidth,ht=2.3ex,dp=0.9ex,leftskip=0.35cm]{foot}\tiny\insertshorttitle\ \textperiodcentered\ \insertshortauthor\end{beamercolorbox}%
    \begin{beamercolorbox}[wd=.50\paperwidth,ht=2.3ex,dp=0.9ex,center]{foot}\tiny\insertsection\end{beamercolorbox}%
    \begin{beamercolorbox}[wd=.16\paperwidth,ht=2.3ex,dp=0.9ex,right,rightskip=0.35cm]{foot}\tiny\insertframenumber\,/\,\inserttotalframenumber\end{beamercolorbox}}}
  \setbeamertemplate{title page}{%
    \vspace*{0.5cm}
    {\usebeamercolor[fg]{title}\usebeamerfont{title}\inserttitle\par}
    \vspace{0.25cm}
    {\usebeamercolor[fg]{subtitle}\large\insertsubtitle\par}
    \vspace{0.3cm}{\color{accentc}\rule{0.92\textwidth}{1.2pt}}\par
    \vspace{0.45cm}
    {\small\insertauthor\par}
    \vspace{0.35cm}
    {\small\color{black!60}\insertdate\par}}
  \makeatletter\define@key{beamerframe}{Title}[true]{}\makeatother   % tuỳ chọn khung "Title" của theme hust
\fi
\setbeamertemplate{caption}[numbered]
\beamertemplatenavigationsymbolsempty
\setbeamertemplate{note page}{%
  \vspace*{4mm}\noindent\begin{minipage}[t]{\textwidth}
    {\small\bfseries\color{accentc}Lời thoại gợi ý}\par\vspace{2mm}
    {\small\insertnote\par}
  \end{minipage}}

%% ---------------------------------------------------------------- lệnh dùng chung
\newcolumntype{Y}{>{\raggedright\arraybackslash}X}
\newcolumntype{P}[1]{>{\raggedright\arraybackslash}p{#1}}
\newcolumntype{Q}[1]{>{\centering\arraybackslash}p{#1}}
\newcolumntype{R}[1]{>{\raggedleft\arraybackslash}p{#1}}
\newcommand{\nhan}[2]{{\setlength{\fboxsep}{1.3pt}\colorbox{#1}{\tiny\bfseries\color{white}\rule[-1.4pt]{0pt}{6.3pt}#2}}}
\newcommand{\nhanS}[2]{{\setlength{\fboxsep}{0.8pt}\colorbox{#1}{\tiny\bfseries\color{white}\rule[-1.1pt]{0pt}{5.6pt}\,#2\,}}}
\newcommand{\OKs}{\nhanS{okc}{THÀNH CÔNG}}
\newcommand{\PARTs}{\nhanS{partc}{CÓ ĐIỀU KIỆN}}
\newcommand{\FAILs}{\nhanS{failc}{KHÔNG HIỆU QUẢ}}
\newcommand{\NAs}{\nhanS{nac}{CHƯA ĐO}}
\newcommand{\DUs}{\nhanS{okc}{ĐỦ}}
\newcommand{\MOTPHANs}{\nhanS{partc}{MỘT PHẦN}}
\newcommand{\TAPCONs}{\nhanS{partc}{TẬP CON}}
\newcommand{\NGHIs}{\nhanS{nac}{NGHI LỖI}}
\newcommand{\OK}{\nhan{okc}{THÀNH CÔNG}}
\newcommand{\PART}{\nhan{partc}{CÓ ĐIỀU KIỆN}}
\newcommand{\FAIL}{\nhan{failc}{KHÔNG HIỆU QUẢ}}
\newcommand{\NA}{\nhan{nac}{CHƯA ĐO}}
\newcommand{\ychinh}[1]{}
\newcommand{\ghichu}[1]{\par\vspace{2pt}{\tiny\color{nac}#1\par}}
\newcommand{\code}[1]{\texttt{#1}}
\newcommand{\effbar}[4]{\fill[#3] (0,#1-0.34) rectangle (#2,#1+0.34);\node[anchor=east,align=right] at (\lx,#1) {#4};}
\newcommand{\tg}[2]{{\setlength{\fboxsep}{1.2pt}\colorbox{#1}{\tiny\bfseries\color{white}\rule[-1.2pt]{0pt}{5.8pt}#2}}}
\newcommand{\hop}[3]{\fcolorbox{#1}{#1!5}{\parbox[t][#2][t]{\dimexpr0.5\linewidth-2\fboxsep-2\fboxrule-3pt\relax}{#3}}}

\title[LLMRec]{Gợi ý tin tức sử dụng các kỹ thuật LLM}
\subtitle{Mười hai hướng tiếp cận LLM trên dữ liệu MIND: hướng nào hiệu quả, hướng nào không}
\author[N.\,N. Hán, N.\,P. Nguyên, L.\,V. Đức]{\texorpdfstring{%
    \begin{tabular}{rl}
        \textbf{Người thực hiện:} & Nguyễn Nam Hán -- 20252095M \\
                                  & Nguyễn Phúc Nguyên -- 20261222M \\
                                  & Lê Văn Đức -- 20261233M \\[0.3em]
        \textbf{GVHD:}            & TS. Ngô Văn Linh
    \end{tabular}}{Nguyễn Nam Hán, Nguyễn Phúc Nguyên, Lê Văn Đức}%
}
\date{2026}

\begin{document}

%% ============================================================================ Bìa
\begin{frame}[noframenumbering,Title]
  \titlepage
  \note{Kính chào thầy cô và các bạn. Nhóm em trình bày đề tài gợi ý tin tức sử dụng các kỹ thuật LLM. Câu hỏi xuyên suốt: trong rất nhiều hướng dùng LLM cho hệ gợi ý những năm 2024--2026, hướng nào thật sự có ích cho tin tức khi được đánh giá chặt chẽ; và từ các hướng có ích đó, có thể xây một hệ thống chạy được trên một GPU phổ thông hay không.}
\par\kern0pt% chiều cao khung tính cả phần sâu của hộp cuối
\end{frame}

%% ============================================================================ Nội dung
\begin{frame}[noframenumbering,shrink]{Nội dung trình bày}
\small
\begin{columns}[T,onlytextwidth]
\column{0.5\textwidth}
\begin{enumerate}
  \setlength{\itemsep}{0.5em}
  \item Bài toán, dữ liệu và câu hỏi nghiên cứu
  \item Cách tiếp cận: quy trình 4 bước, 12 nhóm giả thuyết, giao thức đánh giá
  \item Bức tranh tổng thể: hướng nào thành công, hướng nào không
  \item Chẩn đoán điểm nghẽn: bộ mã hóa (encoder) hay tầng tổng hợp lịch sử
\end{enumerate}
\column{0.5\textwidth}
\begin{enumerate}
  \setcounter{enumi}{4}
  \setlength{\itemsep}{0.5em}
  \item Các hướng thành công và nguyên nhân
  \item Các hướng không thành công và nguyên nhân
  \item Từ thành phần đến hệ thống: kết hợp tuyến tính, GBDT, rò rỉ thời gian
  \item Kết quả, so sánh với các công bố, giới hạn và kết luận
\end{enumerate}
\end{columns}
\vspace{0.6cm}
{\footnotesize\color{nac} Phụ lục: thuật ngữ \textperiodcentered{} câu hỏi thường gặp \textperiodcentered{} các lần chạy và nguồn số liệu \textperiodcentered{} tài liệu tham khảo}
\par\kern0pt% chiều cao khung tính cả phần sâu của hộp cuối
\note{Bài trình bày đi theo đúng trình tự nghiên cứu: đặt bài toán, thử các hướng, xem hướng nào đứng vững, chẩn đoán vì sao, rồi ghép các hướng đứng vững thành hệ thống cuối. Phần 3 là phần quan trọng nhất: một bảng cho biết mỗi hướng H1--H12 thành công hay thất bại.}
\end{frame}

%% ============================================================================ Tóm tắt một trang
\begin{frame}[shrink]{Tóm tắt trong một trang}
\hop{keyc}{2.45cm}{{\footnotesize\bfseries\color{keyc}Bài toán}\par\vspace{2pt}\scriptsize
Dự đoán độc giả sẽ click bài nào trong danh sách ứng viên, dựa trên lịch sử đọc (bộ dữ liệu MIND của Microsoft).
Thách thức lớn nhất: \textbf{82.25\%} click ở tập test rơi vào \textbf{bài mới} chưa từng xuất hiện khi huấn luyện, nên mô hình phải hiểu nội dung bài thay vì nhớ ID.}\hfill
\hop{accentc}{2.45cm}{{\footnotesize\bfseries\color{accentc}Cách làm}\par\vspace{2pt}\scriptsize
Mã hóa văn bản bài báo bằng LLM / encoder \textbf{đóng băng trọng số (frozen)}. Thử nghiệm \textbf{12 nhóm kỹ thuật LLM 2024--2026} (H1--H12) và nhóm dữ liệu đầu vào (F):
sàng lọc trên MIND-small, kiểm chứng trên MIND-large. Mọi so sánh đều bắt cặp (paired), có khoảng tin cậy và chia dữ liệu theo thời gian.}

\vspace{6pt}
\hop{okc}{2.45cm}{{\footnotesize\bfseries\color{okc}Kết quả}\par\vspace{2pt}\scriptsize
\textbf{2} hướng thành công rõ rệt (H8, H10), \textbf{2} hướng hiệu quả có điều kiện (F, H3), \textbf{10} hướng không vượt qua baseline mạnh.
Hệ thống 3 tầng (embedding cố định $\to$ reader $\to$ GBDT): nDCG@10 tăng từ \textbf{0.3787 $\to$ 0.4534} ở chế độ offline; đạt \textbf{0.4749} khi dùng thêm tín hiệu click trong ngày (online).}\hfill
\hop{partc}{2.45cm}{{\footnotesize\bfseries\color{partc}Bài học}\par\vspace{2pt}\scriptsize
(1) Điểm nghẽn nằm ở \textbf{tầng tổng hợp lịch sử đọc}, không phải kích thước bộ mã hóa.
(2) Các kỹ thuật \textbf{nén hoặc viết lại} biểu diễn bằng LLM làm mất chi tiết của bài mới.
(3) Cần đánh giá sát thực tế: tập validation có trọng số theo tỷ lệ bài mới, cross-validation theo khối thời gian.}
\par\kern0pt% chiều cao khung tính cả phần sâu của hộp cuối
\note{Nếu chỉ nhớ một slide thì là slide này. Bài toán khó vì hơn 82\% click rơi vào bài mới. Nhóm giữ bộ mã hóa cố định và thử 12 họ kỹ thuật LLM; chỉ hai họ thành công rõ ràng, hai họ có điều kiện, mười họ không hơn đối chứng mạnh. Ghép các hướng đứng vững cho nDCG@10 tăng từ 0.3787 lên 0.4516; bản mới nhất, thêm đoạn mở rộng do LLM viết cho từng bài, đạt 0.4534 offline và 0.4749 online.}
\end{frame}

%%%%%%%%%%%%%%%%%%%%%%%%%%%%%%%%%%%%%%%%%%%%%%%%%%%%%%%%%%%%%%%%%%%%%%%%%%%%%%%
\section{Bài toán và dữ liệu}
%% ============================================================================ Bài toán
\begin{frame}[shrink]{Bài toán gợi ý tin tức}
\begin{columns}[T,onlytextwidth]
\column{0.56\textwidth}
\footnotesize
Mỗi \textbf{lượt hiển thị} (impression) gồm lịch sử đọc $\mathcal H_u=[d_1,\dots,d_L]$ của độc giả $u$ và danh sách ứng viên $\mathcal C_u$.
Mô hình chấm điểm $s_\theta(u,c)$ cho từng ứng viên để xếp hạng; nhãn là bài báo được click. Ba đặc điểm làm bài toán trở nên thách thức:
\begin{itemize}
  \setlength{\itemsep}{1pt}
  \item \textbf{Bài mới chiếm đa số (cold-start):} phần lớn click rơi vào bài chưa từng có trong tập huấn luyện, khiến lọc cộng tác (collaborative filtering) dựa trên ID hoạt động kém.
  \item \textbf{Giá trị tin tức giảm nhanh theo thời gian:} phân phối dữ liệu dịch chuyển liên tục giữa quá khứ (huấn luyện) và tương lai (test).
  \item \textbf{Ràng buộc độ trễ khi phục vụ (serving):} không thể chạy trực tiếp LLM nhiều tỷ tham số cho từng lượt hiển thị.
\end{itemize}
\column{0.41\textwidth}
\centering
\begin{tikzpicture}[font=\scriptsize,>=Stealth,
  doc/.style={draw=keyc,fill=keyc!6,minimum width=0.62cm,minimum height=0.5cm,inner sep=1pt},
  cand/.style={draw=accentc,fill=accentc!5,minimum width=0.62cm,minimum height=0.5cm,inner sep=1pt}]
\node[anchor=west,font=\scriptsize\bfseries,text=keyc] at (-0.1,2.75) {Lịch sử đọc $\mathcal H_u$};
\foreach \i/\lab in {0/$d_1$,1/$d_2$,2/$d_3$,3/$\cdots$,4/$d_L$} \node[doc] at (0.35+\i*0.72,2.25) {\lab};
\node[anchor=west,font=\scriptsize\bfseries,text=accentc] at (-0.1,1.45) {Ứng viên $\mathcal C_u$};
\foreach \i/\lab in {0/$c_1$,1/$c_2$,2/$c_3$,3/$c_4$,4/$c_5$} \node[cand] (c\i) at (0.35+\i*0.72,0.95) {\lab};
\node[below=1pt of c0,font=\tiny,text=okc] {click};
\node[below=1pt of c3,font=\tiny,text=okc] {click};
\draw[->,thick,nac] (1.8,0.45)--(1.8,-0.05);
\node[draw=nac,rounded corners=2pt,fill=white,align=center,text width=3.3cm] at (1.8,-0.55) {chấm điểm $s_\theta(u,c)$ $\to$ xếp hạng $\mathcal C_u$};
\node[font=\tiny,text=nac] at (1.8,-1.25) {minh họa một lượt hiển thị};
\end{tikzpicture}
\end{columns}
\par\kern0pt% chiều cao khung tính cả phần sâu của hộp cuối
\note{Một lượt hiển thị gồm lịch sử đọc của người dùng và danh sách bài được hiển thị; nhiệm vụ là xếp các bài được click lên đầu. Khó vì đa số bài là bài mới, sở thích trôi theo thời gian, và phải phục vụ nhanh. Vì vậy hướng đi của nhóm là dùng LLM để hiểu nội dung bài, nhưng giữ chi phí lúc phục vụ thấp.}
\end{frame}

%% ============================================================================ Dữ liệu
\begin{frame}[shrink]{Dữ liệu MIND và độ đo đánh giá}
\begin{columns}[T,onlytextwidth]
\column{0.5\textwidth}
\scriptsize
\renewcommand{\arraystretch}{1.18}
\begin{tabularx}{\linewidth}{@{}YR{1.65cm}R{1.65cm}@{}}
\toprule
 & \textbf{MIND-small} & \textbf{MIND-large}\\
\midrule
Bài báo & 65.238 & 104.151\\
Người dùng & -- & 1.000.000\\
Impression huấn luyện & -- & 2.232.748\\
Impression test & 73.152 & 376.471\\
Cặp ứng viên test & -- & 14.085.557\\
Click test vào bài mới & 85.9\% & 82.25\%\\
Mục đích sử dụng & sàng lọc & kiểm chứng\\
\bottomrule
\end{tabularx}
\ghichu{Test = tập \code{*\_dev} (nhãn của tập test chính thức không công khai), chỉ đánh giá một lần ở bước cuối. Validation là ngày cuối của tập huấn luyện.}
\column{0.47\textwidth}
\begin{block}{Ý nghĩa các chỉ số đánh giá}
\scriptsize
\begin{itemize}
  \setlength{\itemsep}{2pt}
  \item \textbf{AUC}: xác suất một bài được click xếp trên một bài không được click trong cùng lượt hiển thị (0.5 là ngẫu nhiên).
  \item \textbf{MRR}: trung bình nghịch đảo thứ hạng của bài được click đầu tiên.
  \item \textbf{nDCG@10} (chỉ số chính): chất lượng 10 vị trí đầu; bài được click xếp càng cao, điểm càng lớn.
  \item \textbf{Cold AUC / Warm AUC}: AUC tính riêng trên các bài được click \emph{chưa từng} / \emph{đã từng} xuất hiện khi huấn luyện.
\end{itemize}
\end{block}
\end{columns}
\par\kern0pt% chiều cao khung tính cả phần sâu của hộp cuối
\note{MIND-small dùng để sàng lọc nhanh nhiều hướng; MIND-large dùng để xác nhận và xây hệ cuối. Chỉ số chính là nDCG@10. Cold AUC rất quan trọng vì hơn 80\% click là bài mới: một kỹ thuật làm tăng nDCG nhưng giảm Cold AUC thường chỉ đang học thuộc bài cũ.}
\end{frame}

%% ============================================================================ Câu hỏi nghiên cứu
\begin{frame}[shrink]{Năm câu hỏi nghiên cứu}
\footnotesize
\renewcommand{\arraystretch}{1.3}
\begin{tabularx}{\textwidth}{@{}lYl@{}}
\toprule
 & \textbf{Câu hỏi} & \textbf{Trả lời ở}\\
\midrule
\textbf{RQ1} & Embedding LLM cố định đứng ở đâu so với NRMS, NAML, Fastformer khi được huấn luyện lại chuẩn xác? & Phần 4\\
\textbf{RQ2} & Điểm nghẽn nằm ở bộ mã hóa (encoder), dữ liệu đầu vào, hay tầng tổng hợp lịch sử đọc? & Phần 4, 5\\
\textbf{RQ3} & Các hướng tiếp cận LLM 2024--2026 (H1--H12) có thực sự hiệu quả khi đánh giá nghiêm ngặt trên bài toán tin tức? & Phần 3, 5, 6\\
\textbf{RQ4} & Ở quy mô MIND-large: nên kết hợp tuyến tính hay dùng GBDT, và làm sao để điểm validation đáng tin cậy? & Phần 7\\
\textbf{RQ5} & Tinh chỉnh (fine-tune) bộ mã hóa có thu hẹp được khoảng cách với các mô hình PLM tinh chỉnh end-to-end không? & Phần 6, 8\\
\bottomrule
\end{tabularx}
\par\kern0pt% chiều cao khung tính cả phần sâu của hộp cuối
\note{Năm câu hỏi này là khung của cả bài. Câu hỏi quan trọng nhất là RQ3: các hướng LLM mới có đứng vững không. RQ1 và RQ2 giải thích vì sao; RQ4 và RQ5 nói về việc xây hệ thống và giới hạn còn lại.}
\end{frame}

%%%%%%%%%%%%%%%%%%%%%%%%%%%%%%%%%%%%%%%%%%%%%%%%%%%%%%%%%%%%%%%%%%%%%%%%%%%%%%%
\section{Cách tiếp cận}
%% ============================================================================ Văn liệu -> giả thuyết
\begin{frame}[shrink]{Quy trình nghiên cứu: bốn bước sàng lọc và tích hợp}
\centering
\begin{tikzpicture}[font=\scriptsize,>=Stealth,
  stp/.style={draw=keyc,fill=keyc!5,rounded corners=2pt,text width=2.9cm,minimum height=2.7cm,align=left,inner sep=4pt,anchor=north west},
  res/.style={draw=okc,fill=okc!6,rounded corners=2pt,text width=2.9cm,minimum height=1.3cm,align=left,inner sep=4pt,anchor=north west}]
\node[stp] (b1) at (0,0) {{\footnotesize\bfseries\color{keyc}1. Dựng baseline chuẩn}\par
  {\tiny\color{nac}MIND-small \textperiodcentered{} lần chạy v7}\par\vspace{3pt}
  Huấn luyện lại NRMS, NAML, Fastformer theo đúng cấu hình bài báo gốc; ablation từng bước (RQ1).};
\node[stp,right=0.35cm of b1.north east,anchor=north west] (b2) {{\footnotesize\bfseries\color{keyc}2. Sàng lọc 12 nhóm}\par
  {\tiny\color{nac}MIND-small \textperiodcentered{} v7, v13}\par\vspace{3pt}
  Đánh giá H1--H12 so với baseline mạnh nhất cùng điều kiện; tính CI bắt cặp; đo phương sai theo seed (RQ2, RQ3).};
\node[stp,right=0.35cm of b2.north east,anchor=north west] (b3) {{\footnotesize\bfseries\color{keyc}3. Kiểm chứng quy mô lớn}\par
  {\tiny\color{nac}MIND-large \textperiodcentered{} v14--v16b}\par\vspace{3pt}
  Chạy lại các hướng hiệu quả trên tập dữ liệu lớn gấp $\sim$14 lần; bổ sung BGE-large và nối (concat) hai encoder.};
\node[stp,right=0.35cm of b3.north east,anchor=north west] (b4) {{\footnotesize\bfseries\color{keyc}4. Tích hợp hệ thống}\par
  {\tiny\color{nac}MIND-large \textperiodcentered{} v15--v17}\par\vspace{3pt}
  Kết hợp tuyến tính, GBDT; xử lý lệch phân phối validation và rò rỉ trong ngày; thử fine-tune encoder (RQ4, RQ5).};
\foreach \a/\b in {b1/b2,b2/b3,b3/b4} \draw[->,thick,keyc] (\a.east)--(\b.west);
\node[res,below=0.22cm of b1.south west,anchor=north west] {\textbf{\color{okc}Kết quả:} NRMS refit đạt 0.4078 là mốc chuẩn (cấu hình cũ V10: 0.3685).};
\node[res,below=0.22cm of b2.south west,anchor=north west] {\textbf{\color{okc}Kết quả:} H8 (LSE) vượt trội; đa số hướng sinh văn bản hoặc nén bằng LLM thất bại.};
\node[res,below=0.22cm of b3.south west,anchor=north west] {\textbf{\color{okc}Kết quả:} H10 vượt LSE khi đủ dữ liệu; F và H3 chỉ hiệu quả có điều kiện.};
\node[res,below=0.22cm of b4.south west,anchor=north west] {\textbf{\color{okc}Kết quả:} đạt 0.4516 nDCG@10 offline, 0.4731 online (bản mới nhất v20c: 0.4534 / 0.4749).};
\end{tikzpicture}
\par\kern0pt% chiều cao khung tính cả phần sâu của hộp cuối
\note{Slide này trả lời câu hỏi ``nhóm đã nghiên cứu như thế nào''. Bước 1: dựng baseline đúng, vì so với baseline yếu thì kỹ thuật nào cũng trông tốt. Bước 2: sàng lọc 12 họ trên MIND-small, mỗi họ so với đối chứng mạnh nhất. Bước 3: mang các họ còn đứng vững lên MIND-large để xác nhận. Bước 4: ghép các họ đứng vững thành hệ thống và sửa hai lỗi đánh giá là validation lệch và rò rỉ trong ngày.}
\end{frame}

%% ============================================================================ Văn liệu -> giả thuyết
\begin{frame}[shrink]{Các hướng nghiên cứu LLM 2024--2026 và giả thuyết tương ứng}
\scriptsize
\renewcommand{\arraystretch}{1.08}
\begin{tabularx}{\textwidth}{@{}P{3.5cm}YP{2.4cm}@{}}
\toprule
\textbf{Hướng tiếp cận} & \textbf{Công trình tiêu biểu} & \textbf{Giả thuyết}\\
\midrule
LLM làm bộ mã hóa (encoder) & LLM2Vec (COLM'24), LLMEmb (AAAI'25), LLM2Rec (KDD'25) & RQ2, H2, H12\\
LLM tóm tắt lịch sử người dùng & EmbSum (RecSys'24) & H1, H5\\
Bổ sung tri thức sinh bởi LLM & KAR (RecSys'24), ONCE (WSDM'24), Yada \& Yamana & H6, F\\
RAG, GraphRAG & GraphRAG (Edge và cs.\ '24), KG-RAG (ACL'25) & H3, H4, H4b\\
Semantic ID & TIGER (NeurIPS'23), LETTER (CIKM'24), RPG (KDD'25) & H7, SID-News\\
LLM làm bộ xếp hạng (reranker) & LLMRank (ECIR'24), LLM4Rerank (WWW'25) & H9\\
Tương tác muộn (Late Interaction) & ColBERT (SIGIR'20) & H8\\
Attention theo ứng viên & CAUM (SIGIR'22) & H10\\
Biểu diễn đa sở thích & ComiRec (KDD'20), MINER (Findings ACL'22) & H11\\
Tinh chỉnh PLM end-to-end & PLM-NR (SIGIR'21), UNBERT (IJCAI'21) & H12\\
Kết hợp ngữ nghĩa + cộng tác & PAD (SIGIR'25) & T\\
\bottomrule
\end{tabularx}
\par\kern0pt% chiều cao khung tính cả phần sâu của hộp cuối
\note{Slide này cho biết mỗi giả thuyết đến từ đâu. Nhóm không tự nghĩ ra các hướng: mỗi H tương ứng một dòng công trình 2024--2026 ở các hội nghị lớn. Cột cuối là mã giả thuyết; phần 3 sẽ cho kết luận của từng mã.}
\end{frame}

%% ============================================================================ Mô hình cơ sở
\begin{frame}[shrink]{Mô hình cơ sở (Baselines) và các bộ mã hóa được khảo sát}
\begin{columns}[T,onlytextwidth]
\column{0.49\textwidth}
\begin{block}{Baseline học từ lịch sử click}
\scriptsize
\begin{itemize}
  \setlength{\itemsep}{1pt}
  \item \textbf{NRMS} (EMNLP 2019): multi-head self-attention cho bài báo và người dùng.
  \item \textbf{NAML} (IJCAI 2019): mạng CNN đa góc nhìn trên tiêu đề, tóm tắt, chuyên mục.
  \item \textbf{Fastformer} (2021): additive attention với độ phức tạp tuyến tính.
  \item Cả ba được \textbf{huấn luyện lại} với embedding từ GloVe-300 theo đúng bài báo gốc; \textbf{NRMS refit} được huấn luyện gộp thêm ngày validation.
  \item \textbf{PLM-NR}, \textbf{UNBERT} (2021): tinh chỉnh PLM end-to-end; trích dẫn kết quả từ bài báo gốc.
\end{itemize}
\end{block}
\column{0.48\textwidth}
\begin{block}{Bộ mã hóa văn bản đóng băng (Frozen Encoders)}
\scriptsize
\begin{itemize}
  \setlength{\itemsep}{1pt}
  \item \textbf{BGE-small / base / large} (384 / 768 / 1024 chiều): encoder hai chiều kiến trúc BERT, lấy vector CLS.
  \item \textbf{Qwen3-Embedding-0.6B} (1024 chiều): kiến trúc LLM decoder, lấy vector ở token cuối.
  \item \textbf{TF-IDF + LSA} (256 chiều): baseline đối chứng dựa trên thống kê từ vựng.
  \item \textbf{Qwen3-1.7B}: LLM sinh văn bản, chỉ dùng cho H6 (sinh hồ sơ) và H9 (chấm điểm).
\end{itemize}
\end{block}
\end{columns}
\par\kern0pt% chiều cao khung tính cả phần sâu của hộp cuối
\note{Điểm khác biệt về phương pháp: nhóm huấn luyện lại NRMS, NAML, Fastformer theo đúng công thức bài báo, nên các so sánh với LLM là công bằng. Bộ mã hóa văn bản luôn giữ cố định; ta chỉ học các tầng nhẹ phía trên.}
\end{frame}

%% ============================================================================ Giao thức
\begin{frame}[shrink]{Giao thức đánh giá và cách đọc kết quả thống kê}
\begin{columns}[T,onlytextwidth]
\column{0.52\textwidth}
\scriptsize
\begin{itemize}
  \setlength{\itemsep}{3pt}
  \item \textbf{Chia theo thời gian:} validation là ngày cuối của tập huấn luyện; tập test (\code{*\_dev}) chỉ đánh giá một lần duy nhất ở cuối. Mọi siêu tham số, epoch, trọng số kết hợp và cấu hình stacker đều chọn trên validation.
  \item \textbf{So sánh bắt cặp (paired) trên từng impression:} báo cáo mức chênh lệch $\Delta$ và khoảng tin cậy (CI) 95\% xấp xỉ chuẩn ($\pm1.96\cdot$SE); Cold/Warm AUC tính SE phân cụm theo impression.
  \item \textbf{Chọn cấu hình tốt nhất của $m$ nhóm trên validation}, kiểm định có hiệu chỉnh Bonferroni: $z_{\mathrm{adj}}=\Phi^{-1}\!\left(1-\frac{0.05}{2m}\right)$.
  \item Báo cáo tách biệt \textbf{Cold AUC / Warm AUC}; các mô hình có huấn luyện được chạy 2--3 seed.
  \item \textbf{73 unit test tự động} trước mỗi lần chạy (kiểm tra công thức độ đo, chống rò rỉ nhãn, giới hạn bộ nhớ); giả thuyết kỳ vọng được ghi nhận trước khi chạy.
\end{itemize}
\column{0.45\textwidth}
\begin{exampleblock}{Cách đọc ``$\Delta$ [CI 95\%]''}
\scriptsize
BGE-large so với Qwen3 (cùng kiến trúc reader):\\
\quad $+0.0102\ [+0.0095, +0.0108]$ \;$\Rightarrow$ CI không chứa 0: \textbf{khác biệt có ý nghĩa thống kê}.\par\vspace{3pt}
kNN trong reader, trên vector concat:\\
\quad $+0.0003\ [-0.0001, +0.0007]$ \;$\Rightarrow$ CI chứa 0: \textbf{khác biệt không có ý nghĩa thống kê (n.s.)}.
\end{exampleblock}
\begin{alertblock}{Lưu ý: phương sai theo seed huấn luyện}
\scriptsize CI chỉ phản ánh sai số lấy mẫu trên tập test. Hai seed của cùng một mô hình có thể lệch \textbf{0.007--0.014} nDCG@10 trên MIND-small, nên mức chênh nhỏ hơn ngưỡng này giữa hai mô hình có huấn luyện chưa đủ để kết luận.
\end{alertblock}
\end{columns}
\par\kern0pt% chiều cao khung tính cả phần sâu của hộp cuối
\note{Giao thức này là lý do nhiều hướng ``trông tốt'' trong văn liệu lại không đứng vững ở đây. Cách đọc đơn giản: nếu khoảng tin cậy không chứa 0 thì khác biệt có ý nghĩa thống kê. Nhưng khoảng tin cậy không tính nhiễu do seed huấn luyện, nên với mô hình đã huấn luyện cần thêm biên an toàn.}
\end{frame}

%% ============================================================================ Kiến trúc
\begin{frame}[shrink]{Kiến trúc hệ thống cuối cùng: ba tầng gắn với các giả thuyết}
\centering
\begin{tikzpicture}[>=Stealth,font=\scriptsize,
  box/.style={draw=keyc,fill=keyc!5,rounded corners=2pt,align=center,inner sep=4pt,minimum height=2.75cm,thick,text width=3.45cm}]
\node[box] (s1) {{\footnotesize\bfseries\color{keyc}Tầng 1: Mã hóa offline}\\[5pt]
  văn bản kèm chuyên mục \tg{partc}{F}\\[2pt]
  Qwen3-Emb $\oplus$ BGE-large\\
  (đóng băng, 2048 chiều)\\[2pt]
  + làm giàu vector bằng kNN \tg{partc}{H3}};
\node[box,right=0.72cm of s1] (s2) {{\footnotesize\bfseries\color{keyc}Tầng 2: Chấm điểm}\\[5pt]
  tương tác muộn LSE \tg{okc}{H8}\\[2pt]
  reader theo ứng viên \tg{okc}{H10}\\[2pt]
  item prior\\(không dùng lịch sử)};
\node[box,right=0.72cm of s2] (s3) {{\footnotesize\bfseries\color{keyc}Tầng 3: GBDT stacker}\\[5pt]
  99 đặc trưng: điểm thành phần, ngữ cảnh slate, độ phủ hiển thị (exposure)\\[2pt]
  chọn mô hình trên validation có trọng số bài mới, CV theo khối thời gian};
\draw[->,thick,keyc] (s1)--node[above=1pt,font=\tiny,text=black,inner sep=0pt]{vector} (s2);
\draw[->,thick,keyc] (s2)--node[above=1pt,font=\tiny,text=black,inner sep=0pt]{điểm} (s3);
\draw[->,thick,accentc] (s3.east)--++(0.95,0) node[midway,above=1pt,font=\tiny\bfseries,text=accentc,inner sep=0pt]{xếp hạng};
\node[below=3pt of s1,font=\tiny,text width=3.5cm,align=center] {mỗi bài chỉ mã hóa một lần: $\approx$4 ms (BGE-large) đến 10 ms (Qwen3) trên 1 GPU T4};
\node[below=3pt of s3,font=\tiny,text width=3.5cm,align=center] {khi phục vụ chỉ tính tích vô hướng và duyệt cây; độ trễ end-to-end \textbf{chưa đo}};
\end{tikzpicture}
\par\kern0pt% chiều cao khung tính cả phần sâu của hộp cuối
\note{Đây là hệ thống cuối, để thầy cô thấy trước đích đến. Tầng 1 mã hóa mỗi bài một lần bằng hai encoder giữ cố định. Tầng 2 tạo nhiều góc nhìn về độ phù hợp giữa người dùng và bài, trong đó quan trọng nhất là LSE và reader theo ứng viên. Tầng 3 là GBDT ghép các điểm này với đặc trưng ngữ cảnh và exposure. Các nhãn màu cho biết mỗi thành phần đến từ giả thuyết nào.}
\end{frame}

%%%%%%%%%%%%%%%%%%%%%%%%%%%%%%%%%%%%%%%%%%%%%%%%%%%%%%%%%%%%%%%%%%%%%%%%%%%%%%%
\section{Bức tranh tổng thể}
%% ============================================================================ Bảng tổng kết
\begin{frame}[shrink]{Bức tranh tổng thể: hướng nào thành công, hướng nào không}
\centering
\begin{tikzpicture}[font=\scriptsize,
  col/.style={rounded corners=2pt,minimum height=4.05cm,anchor=north west,inner sep=4pt,align=left}]
\node[col,draw=okc,fill=okc!5,text width=3.6cm] (c1) at (0,0) {%
  {\footnotesize\bfseries\color{okc}Thành công (2)}\par\vspace{4pt}
  \textbf{H8} \textperiodcentered{} LSE, top-k\par
  {\color{black!70}+0.011 đến +0.014 so với mean pooling; không cần tham số học; nhất quán trên 2 encoder, 2 quy mô}\par\vspace{5pt}
  \textbf{H10} \textperiodcentered{} Reader theo ứng viên\par
  {\color{black!70}+0.0105 so với LSE (MIND-large); thành phần đơn lẻ mạnh nhất. MIND-small: vượt LSE trên Qwen3}};
\node[col,draw=partc,fill=partc!5,text width=3.6cm,right=0.15cm of c1.north east,anchor=north west] (c2) {%
  {\footnotesize\bfseries\color{partc}Có điều kiện (2)}\par\vspace{4pt}
  \textbf{F} \textperiodcentered{} Ghép chuyên mục vào văn bản\par
  {\color{black!70}+0.0186 với Qwen3; không cải thiện trên BGE, không cải thiện bài mới}\par\vspace{5pt}
  \textbf{H3} \textperiodcentered{} Truy hồi kNN\par
  {\color{black!70}+0.0037 khi đưa qua reader, chủ yếu tăng Cold AUC; cộng trực tiếp làm giảm điểm}};
\node[col,draw=failc,fill=failc!5,text width=5.3cm,right=0.15cm of c2.north east,anchor=north west] (c3) {%
  {\footnotesize\bfseries\color{failc}Không hiệu quả (10)}\par\vspace{3pt}
  \textbf{H1} lịch sử dạng văn bản: $-0.0211$\par
  \textbf{H2}$^{\dagger}$ LoRA/DoRA cho encoder: chỉ +0.003\par
  \textbf{H4} GraphRAG: val tăng, test $-0.0066$\par
  \textbf{H4b} LLM tóm tắt cụm: không đổi\par
  \textbf{H5}$^{\dagger}$ User Tower LoRA: hơn LSE, kém reader\par
  \textbf{H6}$^{\dagger}$ hồ sơ do LLM sinh: $-0.0308$\par
  \textbf{H7} Semantic ID: kém mean pooling $-0.031$\par
  \textbf{H9}$^{\dagger}$ LLM chấm điểm: $-0.0112$\par
  \textbf{H11} đa sở thích: thấp hơn LSE\par
  \textbf{H12}$^{\dagger}$ mở khóa encoder: +0.0005 (n.s.)};
\end{tikzpicture}
\par\vspace{3pt}

\par\kern0pt% chiều cao khung tính cả phần sâu của hộp cuối
\note{Đây là câu trả lời trực tiếp cho câu hỏi ``hướng nào thành công''. Thành công rõ chỉ có H8 và H10: cả hai đều thuộc tầng tổng hợp lịch sử. F và H3 có ích trong điều kiện cụ thể. Mười hướng còn lại không hơn đối chứng mạnh; nhiều hướng trong số đó vẫn hơn mean pooling, nhưng mean pooling là đối chứng yếu, nên không tính là thành công. Họ T và SID-News chưa có số đo nên không xếp loại.}
\end{frame}

%% ============================================================================ Bảng tra cứu từng hướng
\begin{frame}[shrink]{Tra cứu nhanh: kết luận cho từng hướng H1--H12}
\scriptsize
\renewcommand{\arraystretch}{0.96}
\setlength{\tabcolsep}{3pt}
\begin{tabularx}{\textwidth}{@{}P{0.9cm}P{3.85cm}YP{2.15cm}@{}}
\toprule
\textbf{Mã} & \textbf{Hướng thử nghiệm} & \textbf{Kết quả chính (so với đối chứng mạnh nhất)} & \textbf{Kết luận}\\
\midrule
H1 & Chuyển lịch sử đọc thành văn bản & small $-0.0211$ đến $-0.0014$; large $-0.0101$ so với mean pooling; thêm MLP học được (v20b): $+0.0080$ so với H1, vẫn $-0.0032$ so với mean pooling & \FAILs\\
H2 & LoRA/DoRA cho encoder$^{\dagger}$ & +0.0016 đến +0.0029; Cold AUC giảm & \FAILs\\
H3 & Làm giàu vector bài bằng kNN & qua reader +0.0037; cộng trực tiếp vào vector bài thì giảm & \PARTs\\
H4 & GraphRAG (tâm cụm đồ thị) & val +0.0199 nhưng test $-0.0066$ & \FAILs\\
H4b & LLM tóm tắt chủ đề mỗi cụm & không khác biệt có ý nghĩa thống kê & \FAILs\\
H5 & User Tower dùng LoRA$^{\dagger}$ & 0.5\% dữ liệu: ngang LSE; 3.5\%: hơn LSE +0.0085, kém reader $-0.0048$; không cải thiện khi kết hợp & \FAILs\\
H6 & LLM sinh hồ sơ, mở rộng bài$^{\dagger}$ & hồ sơ $-0.0308$ / $-0.0278$ (small / large); mở rộng toàn bộ 104.151 bài (v20a): +0.0013 so với LSE, reader không đổi; ghép vào hệ thống (v20c): online +0.0018, offline +0.0004 & \FAILs\\
H7 & Semantic ID + SID-GPT & small: 0.2898, ngang ngẫu nhiên (0.2897); large (v20b, chấm tầng chủ đề): 0.3473, $-0.0314$ so với mean pooling & \FAILs\\
\textbf{H8} & \textbf{Tương tác muộn (LSE)} & \textbf{+0.011 đến +0.014, nhất quán trên 2 encoder và 2 quy mô} & \OKs\\
H9 & LLM chấm điểm ứng viên$^{\dagger}$ & small $-0.0112$; large $-0.0092$ (n.s.); cả hai trên $n=2.000$ & \FAILs\\
\textbf{H10} & \textbf{Reader theo ứng viên} & \textbf{vượt LSE +0.008 đến +0.013 (trừ BGE-small)} & \OKs\\
H11 & Biểu diễn đa sở thích & K-Means $-0.0024$ so với LSE; reader $K=4$ n.s. & \FAILs\\
H12 & Mở khóa 2 tầng cuối BGE-large$^{\dagger}$ & +0.0005 (n.s.); Cold AUC $-0.0069$ & \FAILs\\
F & Thêm chuyên mục vào văn bản & Qwen3 +0.0186, BGE $-0.0024$ so với LSE & \PARTs\\
\bottomrule
\end{tabularx}
\par\kern0pt% chiều cao khung tính cả phần sâu của hộp cuối
\note{Bảng này dùng để tra cứu khi chuẩn bị hoặc khi bị hỏi về một hướng cụ thể. Đọc theo hàng: hướng đã thử, con số quan trọng nhất so với đối chứng mạnh nhất, và kết luận. Chỉ H8 và H10 thành công; F và H3 có điều kiện; T và SID-News chưa đo.}
\end{frame}

%% ============================================================================ Quy mô dữ liệu của các hướng LLM
\begin{frame}[shrink]{Quy mô dữ liệu thực nghiệm của từng hướng LLM}
\scriptsize
\renewcommand{\arraystretch}{1.0}
\setlength{\tabcolsep}{3pt}
\begin{tabularx}{\textwidth}{@{}P{3.75cm}YP{2.55cm}P{1.45cm}@{}}
\toprule
\textbf{Hướng tiếp cận} & \textbf{Dữ liệu huấn luyện đã dùng} & \textbf{Đánh giá trên test} & \textbf{Mức độ}\\
\midrule
Qwen3-0.6B, BGE-large làm encoder (H8, H10, F) & mã hóa toàn bộ 104.151 bài; reader học 1.2 triệu mẫu (bản mở rộng: đủ 2.71 triệu) & toàn bộ 376.471 & \DUs\\
Qwen3-4B làm encoder (RQ2, v18a) & mã hóa toàn bộ 104.151 bài; reader học 1.2 triệu mẫu & toàn bộ & \DUs\\
H1 chuyển lịch sử thành văn bản & zero-shot; bản thêm MLP (v20b): 300.000 lịch sử huấn luyện & small 41\%; large 16\%; v20b: toàn bộ & \DUs\\
H4b LLM tóm tắt cụm & không cần huấn luyện; tóm tắt toàn bộ cụm & toàn bộ (small) & \DUs\\
\midrule
H2 LoRA/DoRA cho encoder & $\le$ 1.500 bước $\times$ 32 = 48.000 cặp đồng click (co-click); chỉ trên small & toàn bộ (small) & \MOTPHANs\\
H5 User Tower LLM + LoRA & v14: 12.800 lịch sử (0.5\%); v19b: 93.952 (3.5\%), 200 phút & v14: 30.000; v19b: toàn bộ & \MOTPHANs\\
H12 mở khóa 2 tầng BGE-large & 576.000 mẫu ($\approx$ 0.5 epoch); truyền ngược gradient phía ứng viên & toàn bộ & \MOTPHANs\\
H6 LLM mở rộng bài, viết hồ sơ & zero-shot; mở rộng toàn bộ 104.151 bài (v20a, 96 phút trên 2 GPU) và tích hợp vào hệ thống (v20c); hồ sơ: chỉ tập con & mở rộng bài: toàn bộ; hồ sơ: 2.500 & \DUs\\
H9 LLM chấm điểm Có/Không & zero-shot; chỉ rerank top-10 của mean pooling & 2.000 (large: 0.5\%) & \TAPCONs\\
H7 Semantic ID + SID-GPT & 3 epoch (small); 4 epoch, 18.392 bước (large, v20b) & toàn bộ (small, large) & \DUs\\
\bottomrule
\end{tabularx}
\ghichu{LLM embedding: Qwen3-Embedding-0.6B (H1, H5); LLM sinh: Qwen3-1.7B (H4b, H6, H9). Tỷ lệ phần trăm tính trên tập test (73.152 impression ở small, 376.471 ở large) hoặc trên 2.71 triệu mẫu huấn luyện. Ước tính thời gian chạy toàn bộ trên GPU T4: H5 (1 epoch) $\approx$ 87 giờ; H9 trên toàn bộ test $\approx$ 38 giờ; mở rộng cả 104.151 bài: đo thật 96 phút trên 2 GPU T4 (v20a). Giới hạn một phiên Kaggle là 12 giờ.}
\par\kern0pt% chiều cao khung tính cả phần sâu của hộp cuối
\note{Slide này để đánh giá công bằng các hướng LLM. Các hướng thành công đều chạy trên toàn bộ dữ liệu. Ngược lại, các hướng huấn luyện LLM mới dùng một phần nhỏ: tháp người dùng mới thấy 3.5\% dữ liệu huấn luyện, mở khóa encoder khoảng nửa epoch. LLM chấm điểm chỉ đo trên vài nghìn impression và đặt trên nền mean pooling; riêng LLM viết mở rộng bài đã chạy đủ toàn bộ 104.151 bài ở v20a và được ghép vào hệ thống ở v20c. Vì vậy với các hướng này, nhóm chỉ kết luận trong ngân sách đã thử; chạy đủ cần nhiều GPU hơn một phiên Kaggle, ví dụ một epoch của tháp người dùng ước khoảng 87 giờ trên T4.}
\end{frame}

%% ============================================================================ Biểu đồ độ lớn
\begin{frame}[shrink]{Mức độ tác động của từng nhóm kỹ thuật (\texorpdfstring{$\Delta$}{Delta} nDCG@10)}
\centering
\begin{tikzpicture}[x=133cm,y=-0.33cm,font=\tiny]
\def\lx{-0.0362}
\foreach \v in {-0.03,-0.02,-0.01,0.01,0.02} \draw[black!12] (\v,0.4)--(\v,14.6);
\draw[black!55] (0,0.4)--(0,14.6);
\effbar{1}{0.0186}{partc!85}{F \textperiodcentered{} chuyên mục, Qwen3 (so với LSE, MIND-large)}
\effbar{2}{0.0113}{okc!85}{H8 \textperiodcentered{} LSE (so với mean pooling, BGE-large)}
\effbar{3}{0.0105}{okc!85}{H10 \textperiodcentered{} reader (so với LSE, BGE-large)}
\effbar{4}{0.0037}{partc!85}{H3 \textperiodcentered{} kNN trong reader (BGE-large)}
\effbar{5}{0.0029}{failc!80}{H2 \textperiodcentered{} LoRA/DoRA (so với mean pooling, MIND-small)}
\effbar{6}{0.0005}{failc!80}{H12 \textperiodcentered{} mở khóa 2 tầng (so với nhánh đóng băng)}
\effbar{7}{-0.0002}{failc!80}{H5 \textperiodcentered{} thêm User Tower LoRA vào bộ kết hợp}
\effbar{8}{-0.0024}{partc!85}{F \textperiodcentered{} chuyên mục, BGE (so với LSE, MIND-large)}
\effbar{9}{-0.0024}{failc!80}{H11 \textperiodcentered{} K-Means $K{=}5$ (so với LSE, BGE-large)}
\effbar{10}{-0.0066}{failc!80}{H4 \textperiodcentered{} GraphRAG (so với mean pooling, MIND-large)}
\effbar{11}{-0.0112}{failc!80}{H9 \textperiodcentered{} LLM chấm điểm (MIND-small, $n{=}2.000$)}
\effbar{12}{-0.0211}{failc!80}{H1 \textperiodcentered{} lịch sử dạng văn bản (MIND-small)}
\effbar{13}{-0.0308}{failc!80}{H6 \textperiodcentered{} hồ sơ do LLM sinh (MIND-small)}
\effbar{14}{-0.0340}{failc!80}{H7 \textperiodcentered{} SID-GPT (so với mean pooling, MIND-small)}
\node[anchor=west] at (0.0188,1) {+0.0186};
\node[anchor=west] at (0.0115,2) {+0.0113};
\node[anchor=west] at (0.0107,3) {+0.0105};
\node[anchor=west] at (0.0039,4) {+0.0037};
\node[anchor=west] at (0.0031,5) {+0.0029};
\node[anchor=west] at (0.0007,6) {+0.0005};
\foreach \y/\lab in {7/$-0.0002$,8/$-0.0024$,9/$-0.0024$,10/$-0.0066$,11/$-0.0112$,12/$-0.0211$,13/$-0.0308$}
  \node[anchor=west] at (0.0006,\y) {\lab};
\node[anchor=west] at (0.0006,14) {$-0.1011$ (thanh bị cắt gọn)};
\draw[white,line width=1.2pt] (-0.0318,13.62)--(-0.0326,14.38);
\draw[white,line width=1.2pt] (-0.0328,13.62)--(-0.0336,14.38);
\foreach \v/\lab in {-0.03/$-0.03$,-0.02/$-0.02$,-0.01/$-0.01$,0/0,0.01/+0.01,0.02/+0.02} \node[below] at (\v,14.6) {\lab};
\end{tikzpicture}
\ghichu{Mỗi thanh được so với đối chứng trực tiếp của chính nó (ghi trong nhãn). Biểu đồ minh họa chiều hướng và biên độ tác động; không dùng để so sánh ngang giữa các nhóm khác đối chứng. Màu: xanh = thành công, vàng = có điều kiện, đỏ = không hiệu quả.}
\par\kern0pt% chiều cao khung tính cả phần sâu của hộp cuối
\note{Biểu đồ cho thấy độ lớn. Các thanh dương lớn nhất đều thuộc tầng tổng hợp lịch sử và dữ liệu đầu vào. Các hướng dùng LLM để sinh hoặc nén biểu diễn nằm ở phía âm; Semantic ID giảm tới 0.1011 nên thanh bị cắt bớt. F xuất hiện hai lần để cho thấy vì sao nó chỉ có điều kiện: dương với Qwen3, âm với BGE.}
\end{frame}

%% ============================================================================ Vì sao thất bại
\begin{frame}[shrink]{Vì sao nhiều hướng tiếp cận LLM không hiệu quả trên MIND?}
\hop{failc}{2.5cm}{{\footnotesize\bfseries\color{failc}1. Nén hoặc viết lại biểu diễn bài báo}\par\vspace{2pt}\scriptsize
GraphRAG thay bài báo bằng tâm cụm, Semantic ID lượng tử hóa thành mã rời rạc, LLM tóm tắt lịch sử và hồ sơ thành đoạn văn.
Khi 82.25\% click rơi vào bài mới, việc làm mất chi tiết đặc trưng của từng bài sẽ làm giảm chất lượng xếp hạng (GraphRAG: Cold AUC $-0.0181$).}\hfill
\hop{partc}{2.5cm}{{\footnotesize\bfseries\color{partc}2. Chỉ vượt qua baseline yếu}\par\vspace{2pt}\scriptsize
K-Means đa sở thích: +0.0090 so với mean pooling nhưng $-0.0024$ so với LSE. User Tower LoRA (0.5\% dữ liệu) vượt mean pooling nhưng chỉ ngang LSE.
Bản thân phép gộp LSE (không tham số) đã thu được phần lợi ích này.}

\vspace{6pt}
\hop{keyc}{2.5cm}{{\footnotesize\bfseries\color{keyc}3. Tập validation ``ấm'' gây sai lệch}\par\vspace{2pt}\scriptsize
Ở lần chạy v14, so với mean pooling: kNN RAG đạt $+0.0130$ trên validation nhưng giảm $-0.0037$ trên test; GraphRAG từ $+0.0199 \to -0.0066$;
reader + kNN từ $+0.0653 \to +0.0302$. Nguyên nhân do validation có tỷ lệ bài mới thấp hơn test rất nhiều.}\hfill
\hop{nac}{2.5cm}{{\footnotesize\bfseries\color{nac}4. Tinh chỉnh với ngân sách tính toán hạn chế}\par\vspace{2pt}\scriptsize
Mở khóa 2 tầng cuối BGE-large ($\approx$ 0.5 epoch): +0.0005 (n.s.), Cold AUC $-0.0069$.
User Tower LoRA: học 0.5\% dữ liệu chỉ ngang LSE, học 3.5\% vượt LSE +0.0085 nhưng vẫn kém reader. Cả hai \textbf{chưa chạy hết dữ liệu}, nên chưa đủ cơ sở bác bỏ bản chất phương pháp.}
\par\kern0pt% chiều cao khung tính cả phần sâu của hộp cuối
\note{Bốn mẫu thất bại lặp lại nhiều lần. Một: nén hoặc viết lại biểu diễn làm mất chi tiết của bài mới. Hai: nhiều hướng chỉ thắng mean pooling, còn LSE không tham số đã làm được điều đó. Ba: validation có ít bài mới nên đánh giá quá cao các mô hình mạnh ở bài cũ. Bốn: tinh chỉnh trong ngân sách GPU nhỏ chưa đủ để thay đổi kết quả.}
\end{frame}

%%%%%%%%%%%%%%%%%%%%%%%%%%%%%%%%%%%%%%%%%%%%%%%%%%%%%%%%%%%%%%%%%%%%%%%%%%%%%%%
\section{Chẩn đoán điểm nghẽn}
%% ============================================================================ RQ1 baseline đầy đủ
\begin{frame}[shrink]{RQ1 \textperiodcentered{} NRMS, NAML, Fastformer huấn luyện lại so với embedding cố định}
\scriptsize
\renewcommand{\arraystretch}{1.04}
\setlength{\tabcolsep}{4pt}
\begin{tabularx}{\textwidth}{@{}YQ{1.1cm}Q{1.0cm}Q{1.32cm}Q{1.35cm}Q{1.05cm}@{}}
\toprule
\textbf{Mô hình (MIND-small)} & \textbf{Huấn luyện?} & \textbf{AUC} & \textbf{nDCG@10} & \textbf{Cold AUC} & \textbf{Nguồn}\\
\midrule
NRMS, seed 0 & có & 0.6392 & 0.3937 & 0.6471 & v7\\
NRMS, seed 1 & có & 0.6501 & 0.4005 & 0.6630 & v7\\
NAML & có & 0.6471 & 0.4006 & 0.6698 & v7\\
Fastformer (trung bình 2 seed) & có & 0.6426 & 0.3927 & 0.6623 & v13\\
NRMS refit (học gộp thêm ngày validation) & có & 0.6555 & 0.4078 & 0.6700 & v7\\
Fastformer refit & có & 0.6550 & 0.4033 & 0.6767 & v13\\
\midrule
BGE-small + LSE (H8) & không & 0.6386 & 0.4037 & 0.6675 & v7\\
Qwen3 + chuyên mục + LSE (F, H8) & không & 0.6556 & \textbf{0.4094} & 0.6637 & v13\\
\bottomrule
\end{tabularx}
\begin{itemize}
  \setlength{\itemsep}{1pt}
  \item LSE vượt NRMS seed 0 (+0.010) và Fastformer (+0.011), tương đương NAML và NRMS seed 1 (+0.003; hai seed NRMS lệch nhau 0.0068); kém NRMS refit $-0.0041\ [-0.0061, -0.0021]$.
  \item Qwen3 + chuyên mục + LSE ngang với NRMS refit (+0.0016, CI chứa 0) và vượt Fastformer refit $+0.0061\ [+0.0022, +0.0100]$.
\end{itemize}
\ghichu{Dùng GloVe-300 theo đúng bài báo gốc. Kết quả bản cũ V10: NRMS 0.3685, NAML 0.3953, Fastformer 0.3642. Số liệu tham khảo từ các bài báo: NRMS 0.410--0.412, NAML 0.412--0.415.}
\par\kern0pt% chiều cao khung tính cả phần sâu của hộp cuối
\note{Đây là câu trả lời đầy đủ cho RQ1. Cả ba baseline NRMS, NAML, Fastformer đều được huấn luyện lại theo công thức bài gốc với GloVe. Embedding cố định cộng LSE, không học tham số nào, ngang các baseline huấn luyện một lần; chỉ NRMS khi được học thêm các ngày validation là còn hơn. Bản zero-shot tốt nhất, Qwen3 cộng chuyên mục cộng LSE, ngang NRMS refit và hơn Fastformer refit có ý nghĩa thống kê.}
\end{frame}

%% ============================================================================ RQ2 encoder sweep
\begin{frame}[shrink]{RQ2 \textperiodcentered{} Tăng kích thước bộ mã hóa không cải thiện được mean pooling}
\begin{columns}[T,onlytextwidth]
\column{0.53\textwidth}
\scriptsize
Giữ cố định phép gộp mean pooling, thay đổi bộ mã hóa (MIND-small):\par\vspace{2pt}
\renewcommand{\arraystretch}{1.18}
\setlength{\tabcolsep}{4pt}
\begin{tabularx}{\linewidth}{@{}YQ{1.35cm}R{1.85cm}@{}}
\toprule
\textbf{Bộ mã hóa (số chiều)} & \textbf{nDCG@10} & \textbf{$\Delta$ so với}\\
 & & \textbf{BGE-small}\\
\midrule
TF-IDF + LSA (256) & 0.3467 & $-0.0443$\\
BGE-base (768) & -- & $-0.0185$\\
Qwen3-Embedding-0.6B (1024) & 0.3784 & $-0.0125$\\
BGE-small (384) & 0.3909 & mốc\\
BGE-large (1024) & 0.3924 & $+0.0015$ (n.s.)\\
\bottomrule
\end{tabularx}
\column{0.44\textwidth}
\footnotesize
\begin{itemize}
  \setlength{\itemsep}{3pt}
  \item Tăng gấp $\sim$10 lần số tham số (BGE-small $\to$ BGE-large) chỉ tăng +0.0015 (n.s.); BGE-base thậm chí thấp hơn BGE-small.
  \item Trên MIND-large, khi làm đầu vào cho reader, BGE-large (335M tham số) vẫn vượt Qwen3-Embedding-0.6B: $+0.0102\ [+0.0095, +0.0108]$.
  \item Trên MIND-large (v18a): \textbf{Qwen3-Embedding-4B} làm đầu vào reader chỉ ngang bản 0.6B ($+0.0006$) và kém BGE-large ($-0.005$ đến $-0.009$).
  \item Khoảng cách tới NRMS refit (0.4078) nằm ở \textbf{tầng tổng hợp lịch sử}: mean pooling cào bằng mọi bài đã đọc, trong khi NRMS học trọng số attention để lọc nhiễu.
\end{itemize}
\end{columns}
\par\kern0pt% chiều cao khung tính cả phần sâu của hộp cuối
\note{Nếu LLM lớn hơn tự động tốt hơn thì Qwen3 0.6B phải hơn BGE-small; thực tế ngược lại. Tăng kích thước BGE lên 10 lần cũng gần như không đổi. Vậy phần đáng cải thiện là cách gộp lịch sử đọc thành điểm số, và đó là chỗ H8, H10 tác động.}
\end{frame}

%%%%%%%%%%%%%%%%%%%%%%%%%%%%%%%%%%%%%%%%%%%%%%%%%%%%%%%%%%%%%%%%%%%%%%%%%%%%%%%
\section{Các hướng thành công}
%% ============================================================================ H8 trực giác
\begin{frame}[shrink]{H8 \textperiodcentered{} Vì sao mean pooling làm loãng các sở thích ngách}
\scriptsize
\[
s_{\mathrm{mean}}(u,c)=\frac{1}{|\mathcal H_u|}\sum_{i\in\mathcal H_u}\langle\mathbf v_i,\mathbf v_c\rangle
\qquad\text{và}\qquad
s_{\mathrm{LSE}}(u,c)=\tau\Big(\log\sum_{i\in\mathcal H_u}\exp\big(\langle\mathbf v_i,\mathbf v_c\rangle/\tau\big)-\log|\mathcal H_u|\Big)
\]
\vspace{-2pt}
\begin{center}
\begin{tikzpicture}[font=\scriptsize,
  pane/.style={rounded corners=2pt,text width=6.3cm,minimum height=1.75cm,align=left,inner sep=5pt,anchor=north west}]
\node[pane,draw=black!45,fill=black!3] (L) at (0,0) {%
  {\bfseries\color{black!75}Mean pooling: tín hiệu bị pha loãng}\par\vspace{3pt}
  Lịch sử: 40 bài thể thao ($\langle\mathbf v_i,\mathbf v_c\rangle\approx0.1$) và 5 bài tài chính ($\approx0.8$).\par
  Ứng viên là bài tài chính: $s_{\mathrm{mean}}=\frac{40\cdot0.1+5\cdot0.8}{45}\approx\mathbf{0.178}$: bị 40 bài thể thao lấn át.};
\node[pane,draw=okc,fill=okc!5,thick] (R) at (6.75,0) {%
  {\bfseries\color{okc}LSE (smooth-max, ví dụ $\tau=0.1$)}\par\vspace{3pt}
  Khuếch đại theo hàm mũ $\exp(\langle\mathbf v_i,\mathbf v_c\rangle/\tau)$ trước khi cộng, giúp vài bài khớp chủ đề chiếm ưu thế.\par
  Cùng ứng viên: $s_{\mathrm{LSE}}\approx\mathbf{0.58}$: giữ được sở thích ngách.};
\end{tikzpicture}
\end{center}
\vspace{-4pt}
\begin{itemize}
  \setlength{\itemsep}{1pt}
  \item Độc giả thường đọc nhiều chủ đề: độ tương đồng cao với một vài bài cùng chủ đề bị chia nhỏ bởi hệ số $1/|\mathcal H_u|$.
  \item LSE nội suy mượt giữa trung bình cộng ($\tau$ lớn) và MaxSim ($\tau\to0$); $\tau$ thực tế được chọn trên tập validation. Phép gộp Top-k (trung bình $k$ bài giống nhất) hoạt động theo cùng nguyên lý.
\end{itemize}
\par\kern0pt% chiều cao khung tính cả phần sâu của hộp cuối
\note{Ví dụ: độc giả đọc 40 bài thể thao và 5 bài tài chính. Với ứng viên tài chính, mean pooling cho điểm thấp vì 40 bài thể thao kéo trung bình xuống. LSE khuếch đại các bài rất giống ứng viên nên giữ được sở thích ngách. Đây là thay đổi không có tham số học.}
\end{frame}

%% ============================================================================ H8 kết quả
\begin{frame}[shrink]{H8 \textperiodcentered{} Kết quả: LSE và top-k cải thiện ổn định, không cần huấn luyện}
\scriptsize
\renewcommand{\arraystretch}{1.05}
\setlength{\tabcolsep}{4pt}
\begin{tabularx}{\textwidth}{@{}P{2.6cm}YQ{1.35cm}Q{1.35cm}P{3.7cm}@{}}
\toprule
\textbf{Toán tử tổng hợp} & \textbf{Ý tưởng} & \textbf{MIND-small} & \textbf{MIND-large} & \textbf{$\Delta$ MIND-large [CI 95\%]}\\
(không tham số) & & BGE-small & BGE-large & so với mean pooling\\
\midrule
\code{frozen} (đối chứng) & trung bình mọi bài đã đọc & 0.3909 & 0.3933 & mốc\\
\code{pool\_recency} & trọng số giảm dần theo thời gian & 0.3917 & 0.3933 & $-0.0000\ [-0.0002, +0.0001]$\\
\code{pool\_maxsim} & chỉ lấy bài giống nhất ($k=1$) & 0.3941 & 0.3916 & $-0.0017\ [-0.0025, -0.0009]$\\
\code{pool\_km5} (H11) & phân cụm lịch sử thành 5 tâm & 0.4007 & 0.4023 & $+0.0090\ [+0.0085, +0.0094]$\\
\textbf{\code{pool\_topk}} & \textbf{trung bình $k$ bài giống nhất} & \textbf{0.4016} & \textbf{0.4045} & $\mathbf{+0.0111}\ [+0.0105, +0.0118]$\\
\textbf{\code{pool\_lse}} & \textbf{smooth-max (LogSumExp)} & \textbf{0.4037} & \textbf{0.4047} & $\mathbf{+0.0113}\ [+0.0107, +0.0119]$\\
\bottomrule
\end{tabularx}
\begin{itemize}
  \setlength{\itemsep}{1pt}
  \item Nhất quán trên Qwen3-Embedding (MIND-large): top-k +0.0143, LSE +0.0130. Cold AUC của LSE cũng tăng rõ (BGE-large: 0.6632 $\to$ 0.6713).
  \item MaxSim cứng ($k=1$) dễ bị nhiễu bởi một lượt click ngẫu nhiên; K-Means luôn nằm giữa mean pooling và LSE (xem H11).
  \item Trên MIND-small, LSE thu hẹp khoảng 3/4 khoảng cách giữa BGE-small mean pooling (0.3909) và NRMS refit (0.4078).
\end{itemize}
\par\kern0pt% chiều cao khung tính cả phần sâu của hộp cuối
\note{Trên cả hai quy mô và hai bộ mã hóa, LSE và top-k đều tăng khoảng 0.011 đến 0.014 nDCG@10 mà không học gì thêm. MaxSim cứng thì giảm vì quá nhạy với một click nhiễu. Từ đây, LSE được dùng làm đối chứng mạnh: một hướng mới phải hơn LSE mới được coi là có ích.}
\end{frame}

%% ============================================================================ H8 so với attention học được
\begin{frame}[shrink]{H8 so với các tầng attention có huấn luyện trên cùng vector cố định}
\scriptsize
\renewcommand{\arraystretch}{1.02}
\setlength{\tabcolsep}{4pt}
\begin{tabularx}{\textwidth}{@{}YQ{1.15cm}Q{1.55cm}Q{1.55cm}Q{1.55cm}Q{1.55cm}@{}}
\toprule
 & & \multicolumn{2}{c}{\textbf{MIND-small}} & \multicolumn{2}{c}{\textbf{MIND-large}}\\
\cmidrule(lr){3-4}\cmidrule(l){5-6}
\textbf{Cách tổng hợp lịch sử} & \textbf{Tham số} & \textbf{BGE-small} & \textbf{Qwen3} & \textbf{BGE-large} & \textbf{Qwen3}\\
\midrule
mean pooling & không & 0.3909 & 0.3784 & 0.3933 & 0.3787\\
\textbf{LSE (H8)} & không & \textbf{0.4037} & 0.3912 & 0.4047 & 0.3916\\
\midrule
reader kiến trúc NRMS & có & 0.3862 & -- & -- & --\\
reader kiến trúc Fastformer & có & 0.3939 & 0.3944 & -- & --\\
reader $K=4$ truy vấn (H11) & có & 0.3944 & \textbf{0.4029} & -- & --\\
reader theo ứng viên (H10) & có & 0.3968 & 0.3996 & 0.4152 & 0.4050\\
reader theo ứng viên + kNN & có & 0.4033 & -- & \textbf{0.4189} & \textbf{0.4095}\\
\bottomrule
\end{tabularx}
\ghichu{Test nDCG@10 trên cùng bảng embedding cố định. Trung bình 2--3 seed (riêng NRMS reader và + kNN trên MIND-small chạy 1 seed). Nguồn: v7, v13, v16b; ``--'': không chạy thực nghiệm.}
\begin{itemize}
  \setlength{\itemsep}{1pt}
  \item \textbf{BGE-small, MIND-small:} LSE vượt mọi biến thể reader có huấn luyện từ +0.007 đến +0.018 (bản reader + kNN chỉ tương đương).
  \item \textbf{Qwen3 hoặc MIND-large:} reader theo ứng viên vượt LSE từ +0.008 đến +0.013.
  \item Kiến trúc attention cụ thể ít ảnh hưởng hơn lựa chọn bộ mã hóa và quy mô dữ liệu huấn luyện.
\end{itemize}
\par\kern0pt% chiều cao khung tính cả phần sâu của hộp cuối
\note{Slide này trả lời câu hỏi: LSE có thắng được attention học được không. Trên cùng vector cố định, nhóm thử reader kiểu NRMS, kiểu Fastformer, nhiều truy vấn và theo ứng viên. Với BGE-small trên MIND-small, LSE thắng tất cả. Nhưng với Qwen3, và khi có nhiều dữ liệu ở MIND-large, reader theo ứng viên thắng LSE. Kiểu attention cụ thể không quan trọng bằng encoder và lượng dữ liệu. Ví dụ trên BGE-small, LSE hơn reader 4 truy vấn $+0.0093$, CI $[+0.0075, +0.0112]$; với Qwen3 trên MIND-small, cả hai seed của reader theo ứng viên đều trên LSE. Hệ thống cuối vì vậy dùng cả hai: reader làm nền, LSE làm thành phần phụ. Reader kiểu NAML chưa được chạy riêng; NAML chỉ được chạy như baseline đầy đủ ở phần RQ1.}
\end{frame}

%% ============================================================================ H10 thiết kế
\begin{frame}[shrink]{H10 \textperiodcentered{} Reader theo ứng viên: dùng ứng viên làm truy vấn đọc lịch sử}
\begin{columns}[T,onlytextwidth]
\column{0.57\textwidth}
\scriptsize
Tầng mạng nhẹ có huấn luyện đặt trên bảng embedding cố định $\mathbf E$:
\begin{itemize}
  \setlength{\itemsep}{2pt}
  \item \textbf{Reader theo ứng viên} (\code{um\_frozen}, \code{um\_knn\_rag}): chiếu vector qua MLP hai tầng, $\mathbf h_i=f(\mathbf E_{\mathrm{hist},i})$, $\mathbf c=f(\mathbf E_{\mathrm{cand}})$, rồi dùng $\mathbf c$ làm truy vấn (query):
  \[
  \alpha_i=\mathrm{softmax}_i\!\Big(\tfrac{\langle\mathbf W_c\mathbf c,\mathbf h_i\rangle}{\sqrt d}\Big),\quad
  s_{\mathrm{CA}}(u,c)=\Big\langle\textstyle\sum_i\alpha_i\mathbf h_i,\ \mathbf c\Big\rangle .
  \]
  \item \textbf{Reader nén một vector} (\code{um\_nrms}, \code{um\_ff}): dùng khối NRMS hoặc Fastformer nén toàn bộ lịch sử thành một vector $\mathbf u$ duy nhất trước khi xét đến ứng viên.
  \item \textbf{Item prior} (\code{um\_prior}): chỉ dựa vào biểu diễn bài báo, dùng làm đặc trưng cho GBDT.
\end{itemize}
Hàm mất mát: cross-entropy với softmax trên bài được click và các mẫu âm (negative) trong cùng lượt hiển thị:
\[
\mathcal L=-\log\frac{\exp s(u,c^+)}{\exp s(u,c^+)+\sum_{c^-}\exp s(u,c^-)} .
\]
\column{0.4\textwidth}
\centering
\begin{tikzpicture}[font=\scriptsize,>=Stealth,
  h/.style={draw=keyc,fill=keyc!6,minimum width=0.7cm,minimum height=0.5cm,inner sep=1pt}]
\node[draw=accentc,fill=accentc!6,minimum width=1.6cm,minimum height=0.55cm] (c) at (1.7,2.6) {ứng viên $\mathbf c$};
\foreach \i/\lab/\w in {0/$\mathbf h_1$/0.4,1/$\mathbf h_2$/1.6,2/$\mathbf h_3$/0.4,3/$\mathbf h_4$/1.1,4/$\mathbf h_5$/0.4}
  {\node[h] (h\i) at (\i*0.85,0.9) {\lab};
   \draw[->,accentc,line width=\w pt] (c.south)--(h\i.north);}
\node[font=\tiny,text=nac,align=center] at (1.7,0.05) {nét càng đậm: trọng số $\alpha_i$ càng lớn\\(bài cùng chủ đề với ứng viên)};
\end{tikzpicture}
\end{columns}
\par\kern0pt% chiều cao khung tính cả phần sâu của hộp cuối
\note{Reader là phiên bản có học của ý tưởng H8. Thay vì công thức cố định, nó học trọng số attention lên từng bài đã đọc, với chính ứng viên là câu hỏi. Có hai kiểu: theo ứng viên, và nén lịch sử thành một vector trước như NRMS. Câu hỏi là học như vậy có hơn LSE không; kết quả phụ thuộc encoder và lượng dữ liệu.}
\end{frame}

%% ============================================================================ H10 kết quả
\begin{frame}[shrink]{H10 \textperiodcentered{} Reader theo ứng viên vượt LSE trên MIND-large và trên Qwen3}
\begin{columns}[T,onlytextwidth]
\column{0.53\textwidth}
\scriptsize
\textbf{Reader theo ứng viên so với LSE trên cùng encoder} (test nDCG@10):\par\vspace{2pt}
\renewcommand{\arraystretch}{1.12}
\setlength{\tabcolsep}{3pt}
\begin{tabularx}{\linewidth}{@{}YQ{1.05cm}Q{1.15cm}R{1.3cm}@{}}
\toprule
\textbf{Thiết lập} & \textbf{LSE} & \textbf{Reader} & \textbf{$\Delta$}\\
\midrule
MIND-small, BGE-small & 0.4037 & 0.3968 & $-0.0069$\\
MIND-small, Qwen3 & 0.3912 & 0.3996 & $+0.0084$\\
MIND-large, BGE-large & 0.4047 & 0.4152 & $\mathbf{+0.0105}$\\
MIND-large, Qwen3 & 0.3916 & 0.4050 & $\mathbf{+0.0134}$\\
MIND-large, vector concat & -- & \textbf{0.4223} & --\\
\bottomrule
\end{tabularx}
\ghichu{Reader: trung bình 3 seed (BGE-small) hoặc 2 seed; độ lệch giữa hai seed từ 0.001--0.013. $\Delta$ của BGE-large là so sánh bắt cặp (v16b); các $\Delta$ khác tính từ hiệu hai giá trị trung bình.}
\column{0.44\textwidth}
\scriptsize
\begin{itemize}
  \setlength{\itemsep}{2pt}
  \item Với BGE-small và lượng dữ liệu ít (MIND-small), reader chưa vượt được LSE. Với Qwen3 hoặc trên MIND-large (dữ liệu gấp $\sim$14 lần), reader vượt rõ LSE.
  \item Khi đặt trên vector concat hai encoder, reader đạt \textbf{0.4223}; biến thể kết hợp kNN (\code{cat:um\_knn\_rag}) trở thành thành phần nòng cốt của hệ thống.
\end{itemize}
\begin{alertblock}{Lưu ý hiện tượng validation ``ấm''}
\scriptsize Ở lần chạy v14 (Qwen3), so với mean pooling: reader + kNN tăng \textbf{+0.0653} trên validation nhưng chỉ tăng \textbf{+0.0302} trên tập test.
\end{alertblock}
\end{columns}
\par\kern0pt% chiều cao khung tính cả phần sâu của hộp cuối
\note{Kết quả của H10 phụ thuộc encoder và lượng dữ liệu. Với BGE-small trên MIND-small, reader chưa hơn LSE. Với Qwen3, reader đã hơn LSE ngay trên MIND-small; và trên MIND-large, với khoảng 14 lần dữ liệu, reader hơn LSE trên cả hai encoder và trở thành thành phần chính của hệ thống. Cần nhớ validation thổi phồng lợi ích của reader, nên mọi kết luận ở đây dựa trên test.}
\end{frame}

%% ============================================================================ Không gian biểu diễn + H3
\begin{frame}[shrink]{Không gian biểu diễn: kết hợp hai encoder; kNN chủ yếu hỗ trợ bài mới}
\begin{columns}[T,onlytextwidth]
\column{0.55\textwidth}
\scriptsize
Test nDCG@10 trên MIND-large, trung bình 2 seed, cùng kiến trúc reader 128 chiều:\par\vspace{2pt}
\renewcommand{\arraystretch}{1.18}
\begin{tabularx}{\linewidth}{@{}Yccc@{}}
\toprule
\textbf{Reader} & \textbf{Qwen3} & \textbf{BGE-large} & \textbf{Concat (2048)}\\
\midrule
theo ứng viên (\code{um\_frozen}) & 0.4050 & 0.4152 & \textbf{0.4223}\\
+ kNN (\code{um\_knn\_rag}, H3) & 0.4095 & 0.4189 & \textbf{0.4226}\\
văn bản kèm chuyên mục (\code{um\_text}) & 0.4092 & -- & --\\
\bottomrule
\end{tabularx}
\ghichu{Vector concat: $\mathbf e=[\mathbf e_{\mathrm{Qwen3}}\,\Vert\,\mathbf e_{\mathrm{BGE}}]/\sqrt2$. kNN: $\mathbf e'_c=\mathrm{Norm}\big(\mathbf e_c+\beta\,\mathrm{Norm}(\sum_{j\in\mathcal N_k(c)}s_{cj}\mathbf e_j)\big)$, chỉ truy hồi các bài xuất hiện trước thời điểm đánh giá. Hai seed của cùng reader lệch 0.001--0.008.}
\column{0.42\textwidth}
\scriptsize
\begin{itemize}
  \setlength{\itemsep}{3pt}
  \item BGE-large vượt Qwen3 $+0.0102\ [+0.0095, +0.0108]$; nối (concat) hai vector vượt BGE-large $+0.0071\ [+0.0066, +0.0076]$: \textbf{sự đa dạng giữa hai không gian biểu diễn} quan trọng hơn kích thước đơn lẻ.
  \item Làm giàu bằng kNN trong reader: +0.0045 (Qwen3), +0.0037 (BGE-large), nhưng trên vector concat chỉ tăng $+0.0003\ [-0.0001, +0.0007]$; dù vậy Cold AUC vẫn tăng +0.0088.
  \item Áp dụng RAG trực tiếp không qua reader (BGE-large, so với mean pooling): thực thể $-0.0105$, kNN $-0.0017$, GraphRAG $-0.0079$.
  \item Huấn luyện reader trên toàn bộ 2.71 triệu mẫu trong 6 epoch (v18b, v19a): vector concat + kNN đạt \textbf{0.4316} ($+0.0091$), BGE-large $+0.0099$; khi đưa vào hệ thống tổng thể chỉ tăng thêm $+0.0006$ (kết hợp tuyến tính) đến $+0.0016$ (GBDT).
\end{itemize}
\end{columns}
\par\kern0pt% chiều cao khung tính cả phần sâu của hộp cuối
\note{Bảng cho thấy ba điều. BGE-large tốt hơn Qwen3 khi làm đầu vào cho reader, và ghép hai encoder còn tốt hơn nữa. kNN giúp trên từng encoder nhưng gần như hết tác dụng trên vector nối, dù vẫn tăng Cold AUC. Còn dùng RAG trực tiếp không qua reader thì đều giảm.}
\end{frame}

%% ============================================================================ F
\begin{frame}[shrink]{F \textperiodcentered{} Đưa chuyên mục vào văn bản: cải thiện Qwen3, không cải thiện BGE}
\scriptsize
Dữ liệu MIND không có toàn văn: tóm tắt (abstract) trung bình 34.9 từ (5.2\% bị trống), tiêu đề 10.7 từ; gồm 18 chuyên mục chính và 270 chuyên mục con.
Mẫu định dạng \code{cta}: \code{"Category: \{cat\}. Subcategory: \{subcat\}. \{title\}. \{abstract\}"}.\par\vspace{1pt}
\renewcommand{\arraystretch}{1.07}
\begin{tabularx}{\textwidth}{@{}Ycccc@{}}
\toprule
\textbf{Biến thể văn bản ($\Delta$ nDCG@10 so với LSE cùng encoder)} & \textbf{Qwen3, small} & \textbf{BGE, small} & \textbf{Qwen3, large} & \textbf{BGE, large}\\
\midrule
chỉ dùng tiêu đề (bỏ tóm tắt) & $-0.0321$ & $-0.0161$ & -- & --\\
+ chuyên mục (\code{cta}) & $\mathbf{+0.0182}$ & $+0.0013$ & $\mathbf{+0.0186}$ & $-0.0024$\\
+ tên thực thể (\code{ctae} so với \code{cta}) & -- & $+0.0033$ & -- & --\\
mở rộng lịch sử 100 click thay vì 50 & $+0.0013$ & $+0.0011$ & -- & --\\
\bottomrule
\end{tabularx}
\begin{itemize}
  \setlength{\itemsep}{1pt}
  \item Qwen3: kết quả validation đồng nhất với test (+0.0129 và +0.0182); BGE: validation ngược chiều với test ($-0.0096$), kém ổn định.
  \item Mức tăng của Qwen3 tập trung ở các bài báo cũ (Warm AUC tăng từ 0.4939 $\to$ 0.6180); Cold AUC gần như không đổi (+0.0005, n.s.); với BGE-large, Cold AUC giảm $-0.0277$.
  \item Phần tóm tắt đóng vai trò rất quan trọng (loại bỏ làm giảm $-0.016$ đến $-0.032$); kéo dài lịch sử đọc hầu như không mang lại khác biệt.
\end{itemize}
\par\kern0pt% chiều cao khung tính cả phần sâu của hộp cuối
\note{Nhóm kiểm tra xem dùng thêm trường dữ liệu có giúp không. Chuyên mục giúp Qwen3 khoảng 0.018 ở cả hai quy mô, nhưng không giúp BGE, và phần tăng nằm ở bài đã biết chứ không ở bài mới. Tóm tắt thì rất quan trọng; lịch sử dài hơn gần như không đổi.}
\end{frame}

%%%%%%%%%%%%%%%%%%%%%%%%%%%%%%%%%%%%%%%%%%%%%%%%%%%%%%%%%%%%%%%%%%%%%%%%%%%%%%%
\section{Các hướng không thành công}
%% ============================================================================ H2 + H12
\begin{frame}[shrink]{H2, H12 \textperiodcentered{} Tinh chỉnh (fine-tune) bộ mã hóa chưa mang lại lợi ích rõ rệt}
\scriptsize
\textbf{H2} (MIND-small): áp dụng LoRA / DoRA cho encoder với hàm mất mát Matryoshka InfoNCE, lấy hard negative từ cùng impression và học trên các cặp bài đồng click:
chỉ đạt \textbf{+0.0016 đến +0.0029} nDCG@10 so với mean pooling, trong khi Cold AUC giảm; gắn thêm head MLP làm giảm $-0.005$.\par\vspace{3pt}
\textbf{H12} (MIND-large, v17): lưu cache trạng thái ẩn tại tầng 22 của BGE-large (fp16 $\approx$ 20.5 GB), mở khóa 2 tầng cuối và huấn luyện đồng thời với reader trong 18.000 bước;
nhánh đối chứng giữ đóng băng hoàn toàn, chạy cùng số bước và cùng thứ tự mẫu.\par\vspace{3pt}
\renewcommand{\arraystretch}{1.07}
\begin{tabularx}{\textwidth}{@{}lYcccc@{}}
\toprule
\textbf{Nhánh} & \textbf{BGE-large} & \textbf{Val (trọng số bài mới)} & \textbf{Test AUC} & \textbf{Test nDCG@10} & \textbf{Cold AUC}\\
\midrule
\code{ft\_base} (đối chứng) & đóng băng hoàn toàn & 0.4476 & 0.6736 & 0.4208 & 0.6883\\
\code{ft\_reader} & mở khóa 2 tầng cuối & 0.4480 & 0.6710 & 0.4213 & 0.6814\\
\bottomrule
\end{tabularx}
\begin{itemize}
  \setlength{\itemsep}{1pt}
  \item $\Delta$ nDCG@10 $=+0.0005\ [-0.0001, +0.0010]$: không có ý nghĩa thống kê; $\Delta$ Cold AUC $=-0.0069\ [-0.0074, -0.0065]$.
  \item Không vượt qua ngưỡng dừng đặt trước (+0.006). Hiện tượng giảm Cold AUC gợi ý encoder có thể đã khớp quá mức vào phân phối bài cũ và giảm khả năng tổng quát hóa trên bài mới.
\end{itemize}
\par\kern0pt% chiều cao khung tính cả phần sâu của hộp cuối
\note{Đây là thí nghiệm gần nhất với cách làm của PLM-NR: mở khóa hai tầng cuối của BGE-large và huấn luyện cùng reader, có nhánh đối chứng đóng băng chạy cùng điều kiện. Kết quả gần như bằng nhau, Cold AUC còn giảm, và không qua ngưỡng +0.006 đã đặt trước. Vì vậy nhóm không đưa tinh chỉnh vào hệ thống cuối.}
\end{frame}

%% ============================================================================ H5
\begin{frame}[shrink]{H5 \textperiodcentered{} User Tower LLM (LoRA): vượt LSE khi học lâu hơn, nhưng vẫn kém reader}
\scriptsize
\textbf{Thiết lập:} Qwen3-Embedding-0.6B đọc chuỗi lịch sử dạng văn bản, tinh chỉnh bằng LoRA (với attention mask kiểu causal, bidirectional hoặc soft mask).
Trên MIND-large, lần chạy v14 dùng 800 bước $\times$ 16 mẫu = 12.800 mẫu (0.5\% tập huấn luyện), đánh giá trên 30.000 impression test.
Lần chạy v19b huấn luyện 200 phút với 5.872 bước $\times$ 16 = 93.952 mẫu (\textbf{3.5\%} của 2.71 triệu mẫu), đánh giá trên \textbf{toàn bộ} 376.471 impression test.\par\vspace{3pt}
\renewcommand{\arraystretch}{1.07}
\setlength{\tabcolsep}{4pt}
\begin{tabularx}{\textwidth}{@{}P{2.8cm}YQ{0.95cm}Q{1.32cm}P{3.65cm}P{1.75cm}@{}}
\toprule
\textbf{Biến thể} & \textbf{Mô tả} & \textbf{AUC} & \textbf{nDCG@10} & \textbf{$\Delta$ nDCG@10 [CI 95\%]} & \textbf{so với}\\
\midrule
\code{ut\_causal} (v14) & User Tower LoRA, 0.5\% dữ liệu & 0.6272 & 0.3807 & $+0.0065\ [+0.0035, +0.0095]$ & mean pooling\\
\code{ut\_causal+meanpool} & User Tower + mean pooling & 0.6384 & 0.3904 & $+0.0162\ [+0.0138, +0.0187]$ & mean pooling\\
 & & & & $-0.0000\ [-0.0027, +0.0026]$ & \textbf{LSE}\\
\code{super\_apex} & bộ kết hợp + User Tower & 0.6743 & 0.4223 & $-0.0002\ [-0.0009, +0.0006]$ & bộ kết hợp\\
\midrule
\code{ut\_causal} (v19b) & User Tower LoRA, 3.5\% dữ liệu & 0.6470 & 0.4002 & $\mathbf{+0.0085}\ [+0.0077, +0.0094]$ & \textbf{LSE}\\
 & & & & $-0.0048\ [-0.0055, -0.0040]$ & reader Qwen3\\
\code{super\_llm} + tower (v19b) & kết hợp thêm User Tower & 0.6863 & 0.4373 & $0$; các dòng GBDT lệch $\pm0.001$ & hệ thống gốc\\
\bottomrule
\end{tabularx}
\begin{itemize}
  \setlength{\itemsep}{1pt}
  \item Trên MIND-small: User Tower + mean pooling $\approx$ LSE của BGE ($-0.0007$, n.s.), kém NRMS refit ($-0.0035$). Chênh lệch giữa các kiểu attention mask nhỏ hơn phương sai giữa các seed (0.010--0.012); biến thể soft mask không tái lập được kết quả (0.4043 $\to$ 0.3956).
  \item Ở v19b: điểm validation tăng 0.3813 $\to$ 0.3899 $\to$ 0.3906 (tại bước 1.938 / 3.934 / 5.872) và bão hòa khi learning rate giảm dần về 0; cùng mô hình LLM nếu không huấn luyện (H1) thấp hơn $-0.0319$.
\end{itemize}
\par\kern0pt% chiều cao khung tính cả phần sâu của hộp cuối
\note{Tháp người dùng LLM là hướng trông hứa hẹn trong văn liệu. Với 800 bước, tức 0.5\% dữ liệu, nó chỉ bằng LSE. Ở lần chạy v19b, huấn luyện 200 phút, khoảng 3.5\% dữ liệu, và chấm trên toàn bộ test, nó đã hơn LSE rõ, nhưng vẫn kém reader trên vector cố định và không thêm gì cho hệ thống. Cần nói rõ: mô hình mới thấy 3.5\% dữ liệu; một epoch đầy đủ ước khoảng 87 giờ trên T4, vượt khả năng một phiên Kaggle. Đây là kết luận trong ngân sách đã thử, chưa phải kết luận về phương pháp.}
\end{frame}

%% ============================================================================ H11
\begin{frame}[shrink]{H11 \textperiodcentered{} Biểu diễn đa sở thích: tách được các vector sở thích nhưng không tăng điểm}
\begin{columns}[T,onlytextwidth]
\column{0.41\textwidth}
\scriptsize
\textbf{K-Means zero-shot} (\code{pool\_km}$K$): phân cụm lịch sử đọc thành $K$ tâm rồi chấm điểm bằng LSE; trường hợp $K=1$ tương ứng mean pooling, còn $K$ đủ lớn hội tụ về chính LSE.\par\vspace{2pt}
Trên MIND-small: $K=2/3/5$ đạt 0.3958 / 0.3989 / 0.4007, luôn nằm giữa mean pooling (0.3909) và LSE (0.4037).
Trên MIND-large: $K=5$ đạt 0.4023, vẫn thấp hơn LSE (0.4047).\par\vspace{3pt}
\textbf{Reader $K=4$ truy vấn} (\code{um\_mi4}): dùng $K$ vector truy vấn học được để trích xuất đa sở thích; bản \code{mi4d} cộng thêm thành phần phạt trực giao $\lambda\cos^2(\mathbf u_k,\mathbf u_l)$ để buộc các vector sở thích tách biệt nhau.
\column{0.56\textwidth}
\scriptsize
\renewcommand{\arraystretch}{1.15}
\setlength{\tabcolsep}{3pt}
\begin{tabularx}{\linewidth}{@{}YQ{1.32cm}Q{1.42cm}Q{1.42cm}@{}}
\toprule
\textbf{BGE-small, 3 seed} & \textbf{nDCG@10} & \textbf{cos} & \textbf{TV}\\
\midrule
\code{um\_frozen} (1 vector) & 0.3968 & -- & --\\
\code{mi1} ($K=1$, đối chứng) & 0.3906 & -- & --\\
\code{mi4} ($K=4$, không phạt) & 0.3944 & 0.85--0.87 & 0.33--0.34\\
\code{mi4d} ($K=4$, có phạt) & 0.3972 & 0.34--0.39 & 0.71--0.72\\
\bottomrule
\end{tabularx}
\ghichu{cos: độ tương đồng cosine trung bình giữa các vector sở thích; TV: khoảng cách Total Variation giữa các phân phối attention.}
\par\vspace{2pt}
Khi không phạt: cả 4 vector sụp đổ về gần như cùng một hướng (cos 0.85--0.87; trên MIND-large là 0.94--0.96). Khi thêm phạt đa dạng: các nhánh tách biệt rõ nhưng điểm xếp hạng không cải thiện (so với 1 vector: BGE +0.0004, n.s.; Qwen3 $-0.0031$).
\end{columns}
\par\kern0pt% chiều cao khung tính cả phần sâu của hộp cuối
\note{Đa sở thích là hướng rất tự nhiên với độc giả đọc nhiều chủ đề. Nhưng K-Means luôn kém LSE, vì LSE chính là trường hợp có rất nhiều cụm. Reader nhiều truy vấn thì sụp đổ về một sở thích; thêm phạt đa dạng thì tách được, nhưng điểm xếp hạng không tăng.}
\end{frame}

%% ============================================================================ Các hướng LLM sinh
\begin{frame}[shrink]{H1, H4, H6, H7, H9 \textperiodcentered{} Sinh văn bản hoặc nén biểu diễn bằng LLM}
\scriptsize
\renewcommand{\arraystretch}{1.02}
\setlength{\tabcolsep}{4pt}
\begin{tabularx}{\textwidth}{@{}P{2.3cm}P{3.55cm}P{4.55cm}Y@{}}
\toprule
\textbf{Nhóm} & \textbf{Cách làm} & \textbf{Kết quả ($\Delta$ nDCG@10)} & \textbf{Nguyên nhân khả dĩ}\\
\midrule
H1: lịch sử dạng văn bản & nối các tiêu đề đã đọc thành một chuỗi truy vấn dài & small $-0.0211$ đến $-0.0014$ ($n=30.000$); large $-0.0101$ ($n=60.000$); + MLP học được (v20b, toàn bộ test): $-0.0032$ so với mean pooling & chuỗi văn bản dài làm loãng tín hiệu của từng bài riêng lẻ\\
H6: hồ sơ độc giả & dùng Qwen3-1.7B sinh đoạn văn mô tả sở thích & small $-0.0308$; large $-0.0278$ (cả hai $n=2.500$) & tóm tắt văn bản làm mất các sở thích ngách\\
H6: mở rộng bài báo & dùng LLM sinh thêm đoạn mô tả cho mỗi bài & small $+0.0079\ [+0.0021, +0.0137]$; large $+0.0043$ (n.s.); $n=2.500$; toàn bộ (v20a): $+0.0013$ so với LSE, Cold AUC $-0.012$; trong hệ thống (v20c): online $+0.0018\ [+0.0015, +0.0022]$ & thông tin sinh thêm trùng lặp với tín hiệu độ phổ biến; chỉ có ích như thành phần phụ\\
H9: LLM chấm điểm & dùng Qwen3-1.7B chấm xác suất Yes/No cho ứng viên & small $-0.0112$; large $-0.0092$ (n.s.); cả hai $n=2.000$, rerank top-10 của mean pooling & xác suất sinh chưa được hiệu chuẩn (uncalibrated); tốc độ chậm\\
H7: Semantic ID & lượng tử hóa RQ-KMeans; xếp hạng sinh bằng SID-GPT & hồ sơ SID $+0.0035$; SID-GPT 0.2898 $\approx$ ngẫu nhiên 0.2897; large (v20b): chấm tầng chủ đề $+0.0483$ so với cộng cả 3 tầng, vẫn $-0.0314$ so với mean pooling & cách chấm cũ cộng cả tầng mịn của bài mới; mã rời rạc làm mất chi tiết của từng bài\\
H4: GraphRAG & thay vector bài bằng tâm cụm trên đồ thị kNN & small $+0.0043$; large: val $+0.0199$, test $-0.0066$ & lấy tâm cụm làm mờ chi tiết đặc trưng của bài mới\\
H4b: tóm tắt cụm & dùng LLM tóm tắt chủ đề cho mỗi cụm & không khác biệt có ý nghĩa thống kê & --\\
\bottomrule
\end{tabularx}
\par\kern0pt% chiều cao khung tính cả phần sâu của hộp cuối
\note{Nhóm các hướng dùng LLM để viết lại hoặc nén biểu diễn đều không đứng vững. Riêng Semantic ID, trên MIND-small mô hình sinh cho kết quả ngang vector ngẫu nhiên; lần chạy v20b trên MIND-large cho thấy nguyên nhân là cách chấm cộng cả các tầng mã mịn: chỉ chấm tầng chủ đề tăng 0.048, nhưng vẫn kém mean pooling. Mở rộng bài bằng LLM chạy đủ toàn bộ bài ở v20a: riêng lẻ gần như không tăng, ghép vào hệ thống chỉ tăng 0.0018 ở chế độ online.}
\end{frame}



%%%%%%%%%%%%%%%%%%%%%%%%%%%%%%%%%%%%%%%%%%%%%%%%%%%%%%%%%%%%%%%%%%%%%%%%%%%%%%%
\section{Từ thành phần đến hệ thống}
%% ============================================================================ Validation lệch
\begin{frame}[shrink]{Lệch phân phối giữa Validation và Test: 18\% so với 82\% click vào bài mới}
\centering
\begin{tikzpicture}[font=\scriptsize,scale=0.9,transform shape]
\node[anchor=east,font=\scriptsize\bfseries] at (0,0.65) {Validation:};
\fill[keyc] (0.2,0.45) rectangle (1.65,0.85);
\node[white,font=\tiny\bfseries] at (0.92,0.65) {18.07\%};
\fill[black!18] (1.65,0.45) rectangle (8.2,0.85);
\node[black!75] at (4.9,0.65) {81.93\% click rơi vào bài đã xuất hiện trong tập huấn luyện};
\node[anchor=east,font=\scriptsize\bfseries] at (0,0) {Test:};
\fill[accentc] (0.2,-0.2) rectangle (6.78,0.2);
\node[white,font=\scriptsize\bfseries] at (3.5,0) {82.25\% click rơi vào bài mới};
\fill[black!18] (6.78,-0.2) rectangle (8.2,0.2);
\node[black!75,font=\tiny] at (7.49,0) {17.75\%};
\end{tikzpicture}
\par\vspace{4pt}
\begin{columns}[T,onlytextwidth]
\column{0.52\textwidth}
\scriptsize
Mức chênh lệch so với mean pooling trên validation và trên test (lần chạy v14, Qwen3):\par\vspace{2pt}
\renewcommand{\arraystretch}{1.12}
\begin{tabularx}{\linewidth}{@{}Ycc@{}}
\toprule
\textbf{Thành phần} & \textbf{Validation} & \textbf{Test}\\
\midrule
LSE (H8) & $+0.0232$ & $+0.0146$\\
văn bản kèm chuyên mục (F) & $+0.0370$ & $+0.0315$\\
kNN RAG cộng trực tiếp (H3) & $+0.0130$ & $\mathbf{-0.0037}$\\
GraphRAG (H4) & $+0.0199$ & $\mathbf{-0.0066}$\\
reader theo ứng viên (H10) & $+0.0596$ & $+0.0192$\\
reader + kNN & $+0.0653$ & $+0.0302$\\
\bottomrule
\end{tabularx}
\column{0.45\textwidth}
\footnotesize
\begin{itemize}
  \setlength{\itemsep}{3pt}
  \item Tập validation là ngày cuối của cửa sổ huấn luyện nên mang tính ``ấm'': đa số click rơi vào các bài báo đã xuất hiện.
  \item Các thành phần không huấn luyện giữ được phần lớn mức tăng trên test; các mô hình reader mất hơn một nửa biên độ; riêng RAG và GraphRAG \textbf{đảo chiều từ dương sang âm}.
  \item Việc chọn siêu tham số theo trung bình validation thông thường sẽ thiên vị các mô hình chỉ mạnh trên bài cũ.
\end{itemize}
\end{columns}
\par\kern0pt% chiều cao khung tính cả phần sâu của hộp cuối
\note{Đây là phát hiện quan trọng về đánh giá. Validation có chưa tới 20\% click vào bài mới, test có hơn 80\%. Hệ quả là các thành phần mạnh ở bài cũ, như reader hay RAG, trông tốt trên validation hơn hẳn trên test; RAG và GraphRAG còn đổi dấu. Nếu chọn mô hình theo validation thông thường, ta sẽ chọn sai.}
\end{frame}


%% ============================================================================ Hợp nhất tuyến tính
\begin{frame}[shrink]{Kết hợp tuyến tính: một reader cùng ba thành phần zero-shot là đủ}
\scriptsize
Điểm của mỗi thành phần được chuẩn hóa z-score theo từng impression; trọng số được tối ưu bằng coordinate ascent trên tập validation có trọng số bài mới.\par\vspace{3pt}
\renewcommand{\arraystretch}{1.05}
\setlength{\tabcolsep}{3pt}
\begin{tabularx}{\textwidth}{@{}P{2.2cm}YQ{0.95cm}Q{0.95cm}Q{1.32cm}Q{1.05cm}@{}}
\toprule
\textbf{Cấp độ} & \textbf{Thành phần được chọn (trọng số)} & \textbf{Val} & \textbf{AUC} & \textbf{nDCG@10} & \textbf{Cold AUC}\\
\midrule
\code{super\_base} & \code{cat:um\_knn\_rag}, trung bình z-score của 2 seed & 0.4573 & 0.6768 & 0.4260 & 0.7104\\
\textbf{\code{super\_zs}} & + \code{bge:hist100\_lse} 0.1, \code{bge:pool\_lse} 0.1,\newline\phantom{+} \code{qwen3:txt\_cta\_lse} 0.25 & 0.4676 & 0.6863 & \textbf{0.4373} & 0.7185\\
\code{super\_llm}$^{*}$ & + \code{bge:ft\_reader} 0.1 (H12, lần chạy v17) & 0.4679 & 0.6869 & 0.4378 & 0.7181\\
\code{super\_llm+pop} & \code{super\_zs} + popularity online 1.0 & 0.4927 & 0.7133 & 0.4638 & 0.7539\\
\bottomrule
\end{tabularx}
\begin{itemize}
  \setlength{\itemsep}{1pt}
  \item \code{super\_base} vượt điểm của từng seed reader đơn lẻ (0.4226) nhờ lấy trung bình hai seed (+0.0034).
  \item Bổ sung ba thành phần zero-shot giúp tăng $+0.0114\ [+0.0109, +0.0118]$; các reader khác đều nhận trọng số bằng 0; \code{ft\_reader} (H12) chỉ đóng góp thêm +0.0005.
  \item Nếu bỏ nhóm zero-shot khỏi \code{super\_llm+pop}: giảm $-0.0092$; nếu bỏ popularity: giảm $-0.0265$.
\end{itemize}
\ghichu{Val: điểm trên tập validation có trọng số bài mới. (*) Số liệu từ lần chạy v17; các dòng còn lại từ v16b (trong mã nguồn v16b, \code{super\_zs} được đặt tên là \code{super\_llm}).}
\par\kern0pt% chiều cao khung tính cả phần sâu của hộp cuối
\note{Bộ hợp nhất tuyến tính rất đơn giản: chuẩn hóa điểm của từng thành phần rồi cộng có trọng số. Kết quả quan trọng: một reader cộng ba thành phần zero-shot khác loại là đủ, đạt 0.4373; thêm reader cùng loại, kể cả encoder đã tinh chỉnh của H12, gần như không đổi. Đây là mức tốt nhất khi chỉ dùng nội dung và lịch sử.}
\end{frame}

%% ============================================================================ GBDT
\begin{frame}[shrink]{GBDT stacker: 99 đặc trưng thuộc 4 nhóm thông tin}
\footnotesize
\textbf{Vì sao cần GBDT phía trên các reader và bộ kết hợp tuyến tính?} Kết hợp tuyến tính áp dụng chung một bộ trọng số cố định cho mọi ứng viên; trong khi GBDT có thể điều chỉnh trọng số linh hoạt theo ngữ cảnh,
chẳng hạn dựa vào mức độ hiển thị (exposure) của từng bài. Các mô hình reader không bị thay thế mà đóng vai trò là nhóm đặc trưng quan trọng nhất.\par\vspace{3pt}
\scriptsize
\renewcommand{\arraystretch}{1.18}
\begin{tabularx}{\textwidth}{@{}P{3.3cm}cY@{}}
\toprule
\textbf{Nhóm} & \textbf{Số cột} & \textbf{Nội dung}\\
\midrule
\code{scores} (điểm thành phần) & 59 & điểm chuẩn hóa z-score trong từng impression và điểm thô của 41 mô hình thành phần, chủ yếu là các reader trên Qwen3, BGE-large và vector concat\\
\code{ctx} (ngữ cảnh slate) & 15 & độ tương đồng với click gần nhất (\code{last1}, \code{last3}), tỷ lệ trùng khớp chuyên mục, số lượng ứng viên trong slate (\code{ctx\_n\_cand}), mức độ trùng thực thể Wikidata có trọng số IDF\\
\code{meta} (exposure, offline) & 16 & thống kê từ tập huấn luyện (\code{m\_ctr}, \code{m\_unseen}) và động thái hiển thị không dùng nhãn click: tuổi bài báo (\code{m\_age\_h}), tần suất hiển thị trong 1, 6, 24 giờ trước đó\\
\code{pop} (chỉ dùng online) & 9 & log-CTR tích lũy và CTR trượt theo cửa sổ 1, 6, 24 giờ, tính từ các lượt click xảy ra trước đó trong chính ngày test\\
\bottomrule
\end{tabularx}
\par\kern0pt% chiều cao khung tính cả phần sâu của hộp cuối
\note{GBDT nhận 99 đặc trưng chia làm 4 nhóm. Nhóm meta là thông tin hiển thị của bài lấy từ log mà không dùng nhãn, nên vẫn thuộc chế độ offline. Nhóm pop dùng click của chính ngày test nên chỉ thuộc chế độ online. Ba biến thể tương ứng ba mức thông tin.}
\end{frame}

%% ============================================================================ Rò rỉ
\begin{frame}[shrink]{Rò rỉ thông tin trong ngày: vì sao chia fold theo người dùng gây sai lệch}
\scriptsize
GBDT được huấn luyện trên tập validation (150.000 impression $\approx$ 6.05 triệu dòng ứng viên). Việc chia 4 fold theo \textbf{người dùng (user-split)} dẫn đến hiện tượng quá khớp (overfit) vào chính ngày validation đó:
\begin{itemize}
  \setlength{\itemsep}{2pt}
  \item \textbf{Cơ chế gây rò rỉ:} các impression diễn ra cùng thời điểm xuất hiện ở cả tập huấn luyện và tập đánh giá của mỗi fold. Cây quyết định sâu ($d=7$) tận dụng các đặc trưng có độ phân giải cao (tuổi bài, tần suất hiển thị, kích thước slate)
  để nhận diện từng bài báo cụ thể và học thuộc tỷ lệ click của bài đó trong ngày.
  \item \textbf{Biểu hiện:} điểm out-of-fold trên validation cao bất thường (\textbf{0.4922}), nhưng rơi xuống \textbf{0.3738} khi chấm riêng trên các bài chưa xuất hiện ở phần huấn luyện.
\end{itemize}
\begin{block}{Giải pháp khắc phục: Cross-validation theo khối thời gian (Time-block CV)}
\scriptsize
\begin{enumerate}
  \setlength{\itemsep}{1pt}
  \item Chia ngày validation thành \textbf{4 khối thời gian liên tiếp}; loại bỏ một khoảng đệm (tối đa 2\% thời lượng ngày, $\approx$ 29 phút) ở hai biên của khối đánh giá.
  \item Làm thô (quantize) các đặc trưng có nguy cơ định danh bài báo; cấu hình \code{clean} loại bỏ thêm các đặc trưng mang giá trị hằng số trong cùng một impression.
  \item Dừng sớm (early stopping) và chọn cấu hình dựa trên \textbf{forward holdout} (60\% đầu ngày để huấn luyện, 10\% để dừng sớm, 30\% cuối ngày để đánh giá), tối ưu trực tiếp trên \textbf{lượt click vào bài mới}.
  Cấu hình cuối cùng lấy trung bình dự đoán của hai thiết lập \code{v15} và \code{clean}.
\end{enumerate}
\end{block}
\par\kern0pt% chiều cao khung tính cả phần sâu của hộp cuối
\note{Lỗi này khó thấy: fold theo người dùng trông hợp lý nhưng các fold cùng một ngày, cùng các bài, nên cây học thuộc độ phổ biến của từng bài trong ngày đó. Dấu hiệu là điểm validation 0.49 nhưng chỉ 0.37 trên bài mới. Cách sửa là chia fold theo khối thời gian và đánh giá trên bài mới.}
\end{frame}

%% ============================================================================ Rò rỉ kết quả
\begin{frame}[shrink]{So sánh User-split và Time-block CV trên MIND-large}
\scriptsize
\renewcommand{\arraystretch}{1.12}
\setlength{\tabcolsep}{4pt}
\begin{tabularx}{\textwidth}{@{}P{2.95cm}YQ{1.1cm}Q{1.45cm}Q{1.1cm}Q{1.3cm}@{}}
\toprule
\textbf{Biến thể} & \textbf{Cách chia fold} & \textbf{Val} & \textbf{Val bài mới} & \textbf{Test} & \textbf{Test $-$ Val}\\
\midrule
\code{super\_zs} & không dùng GBDT & 0.4275 & -- & 0.4373 & $+0.0098$\\
\code{stack\_meta\_v15cfg} & chia theo người dùng, cây sâu $d=7$ & \textbf{0.4922} & \textbf{0.3738} & 0.4543 & $\mathbf{-0.0379}$\\
\code{stack\_llm} & chia theo khối thời gian & 0.4456 & 0.4109 & 0.4369 & $-0.0087$\\
\code{stack\_meta} & chia theo khối thời gian & 0.4588 & 0.4533 & 0.4516 & $-0.0072$\\
\code{stack\_pop} & chia theo khối thời gian & 0.4682 & 0.4605 & 0.4731 & $+0.0049$\\
\bottomrule
\end{tabularx}
\par\vspace{3pt}
\centering
\begin{tikzpicture}[font=\tiny,x=1.15cm,y=0.5cm]
\node[anchor=east,font=\scriptsize\bfseries,align=right] at (-0.12,2.15) {theo người dùng};
\fill[black!35] (0,2.62) rectangle (4.922,2.94); \node[anchor=west] at (4.97,2.78) {Val \textbf{0.4922} (quá khớp trong ngày)};
\fill[accentc]  (0,2.0) rectangle (3.738,2.32);  \node[anchor=west,text=accentc] at (3.79,2.16) {Val trên bài mới \textbf{0.3738}};
\fill[keyc]     (0,1.38) rectangle (4.543,1.70); \node[anchor=west,text=keyc] at (4.59,1.54) {Test 0.4543 (lệch $-0.0379$)};
\draw[black!25] (-2.3,1.12)--(8.6,1.12);
\node[anchor=east,font=\scriptsize\bfseries,align=right] at (-0.12,0.24) {theo khối thời gian};
\fill[black!35] (0,0.70) rectangle (4.588,1.02); \node[anchor=west] at (4.64,0.86) {Val \textbf{0.4588}};
\fill[accentc]  (0,0.08) rectangle (4.533,0.40); \node[anchor=west,text=accentc] at (4.58,0.24) {Val trên bài mới \textbf{0.4533}};
\fill[keyc]     (0,-0.54) rectangle (4.516,-0.22); \node[anchor=west,text=keyc] at (4.57,-0.38) {Test \textbf{0.4516} (lệch $-0.0072$)};
\end{tikzpicture}
\ghichu{Độ đo: nDCG@10. Val: điểm out-of-fold trên tập validation (với \code{super\_zs} là điểm validation thông thường). Trục ngang biểu đồ bắt đầu từ 0. Toàn bộ số liệu từ lần chạy v16b.}
\par\kern0pt% chiều cao khung tính cả phần sâu của hộp cuối
\note{Bảng xác nhận giả thuyết rò rỉ: fold theo người dùng cho điểm validation 0.4922 nhưng test chỉ 0.4543. Fold theo khối thời gian cho validation gần test. Cần nói thẳng: mô hình mới không cao hơn trên test; điều nhóm được là một con số validation mà mình tin được.}
\end{frame}

%%%%%%%%%%%%%%%%%%%%%%%%%%%%%%%%%%%%%%%%%%%%%%%%%%%%%%%%%%%%%%%%%%%%%%%%%%%%%%%
\section{Kết quả và so sánh}
%% ============================================================================ Kết quả offline
\begin{frame}[shrink]{Kết quả chế độ Offline trên MIND-large (376.471 impression test)}
\scriptsize
\renewcommand{\arraystretch}{1.05}
\setlength{\tabcolsep}{3pt}
\begin{tabularx}{\textwidth}{@{}P{3.05cm}YQ{0.92cm}Q{0.92cm}Q{0.92cm}Q{1.17cm}Q{1.32cm}Q{1.0cm}@{}}
\toprule
\textbf{Mô hình} & \textbf{Nhóm / vai trò} & \textbf{Val} & \textbf{AUC} & \textbf{MRR} & \textbf{nDCG@5} & \textbf{nDCG@10} & \textbf{Cold AUC}\\
\midrule
1. \code{qwen3:frozen} & mốc cơ sở (mean pooling) & 0.3392 & 0.6215 & 0.2922 & 0.3168 & 0.3787 & 0.6583\\
2. \code{bge:frozen} & mean pooling, BGE & 0.3546 & 0.6297 & 0.3042 & 0.3321 & 0.3933 & 0.6632\\
3. \code{cat:frozen} & mean pooling, concat & 0.3560 & 0.6379 & 0.3043 & 0.3348 & 0.3945 & 0.6772\\
\midrule
4. \code{bge:pool\_lse} & H8: LSE & 0.3743 & 0.6459 & 0.3136 & 0.3441 & 0.4047 & 0.6713\\
5. \code{qwen3:txt\_cta\_lse} & F: chuyên mục + LSE & 0.3772 & 0.6554 & 0.3206 & 0.3509 & 0.4102 & 0.6620\\
6. \code{cat:um\_knn\_rag} & H10 + H3: reader & 0.4159 & 0.6734 & 0.3255 & 0.3598 & 0.4226 & 0.7065\\
\midrule
7. \code{super\_zs} & kết hợp tuyến tính & 0.4275 & 0.6863 & -- & -- & 0.4373 & 0.7185\\
8. \code{stack\_llm} & GBDT: điểm + ngữ cảnh & 0.4456 & 0.6853 & 0.3375 & 0.3756 & 0.4369 & 0.7063\\
\textbf{9. \code{stack\_meta}} & \textbf{GBDT + exposure} & \textbf{0.4588} & \textbf{0.7013} & \textbf{0.3503} & \textbf{0.3899} & \textbf{0.4516} & \textbf{0.7248}\\
10. \code{stack\_meta\_v15cfg} & GBDT, chia fold theo user & 0.4922 & 0.7029 & 0.3528 & 0.3930 & 0.4543 & 0.7285\\
\midrule
11. \code{super\_llm} (v20c) & mới nhất: kết hợp tuyến tính & 0.4341 & 0.6939 & 0.3418 & 0.3795 & 0.4407 & 0.7097\\
\textbf{12. \code{stack\_meta} (v20c)} & \textbf{mới nhất: + đoạn LLM viết} & 0.4636 & \textbf{0.7038} & \textbf{0.3504} & \textbf{0.3914} & \textbf{0.4534} & \textbf{0.7256}\\
\bottomrule
\end{tabularx}
\ghichu{Số liệu lấy từ lần chạy v16b (dòng 7 log không xuất MRR và nDCG@5). Val: nDCG@10 trên tập validation (out-of-fold đối với GBDT). Dòng 10 chỉ mang tính tham khảo do điểm validation bị phóng đại bởi rò rỉ thông tin trong ngày. Dòng 11--12: lần chạy v20c (bản mới nhất); dòng 11 dùng quy tắc chọn trên validation không trọng số (đăng ký trước); \code{stack\_meta} v20c hơn hệ tốt nhất trước (v19c) $+0.0004\ [+0.0000, +0.0007]$, 1 seed stacker.}
\par\kern0pt% chiều cao khung tính cả phần sâu của hộp cuối
\note{Bảng đọc từ trên xuống theo đúng các tầng: đối chứng, các thành phần đứng vững, hợp nhất, rồi GBDT. Mô hình offline chính là stack\_meta với 0.4516; bản mới nhất v20c, thêm thành phần có đoạn LLM viết, đạt 0.4534. Dòng 10 cao hơn một chút trên test nhưng điểm validation của nó không đáng tin, nên chỉ để tham khảo.}
\end{frame}

%% ============================================================================ Online
\begin{frame}[shrink]{Chế độ Online: bổ sung tín hiệu phản hồi click trong ngày}
\begin{columns}[T,onlytextwidth]
\column{0.56\textwidth}
\scriptsize
\renewcommand{\arraystretch}{1.18}
\setlength{\tabcolsep}{4pt}
\begin{tabularx}{\linewidth}{@{}Ycccc@{}}
\toprule
\textbf{Mô hình (online)} & \textbf{AUC} & \textbf{MRR} & \textbf{nDCG@10} & \textbf{Cold AUC}\\
\midrule
\code{popularity} (chỉ đếm click, không dùng nội dung) & 0.6540 & 0.3126 & 0.4040 & 0.6986\\
\code{super\_llm+pop} & 0.7133 & 0.3604 & 0.4638 & 0.7539\\
\code{stack\_pop} & 0.7196 & 0.3671 & 0.4731 & 0.7526\\
\textbf{\code{stack\_pop}} (v20c, mới nhất) & \textbf{0.7216} & \textbf{0.3676} & \textbf{0.4749} & 0.7528\\
\bottomrule
\end{tabularx}
\ghichu{So sánh bắt cặp: \code{stack\_pop} $-$ \code{stack\_meta} $=+0.0215\ [+0.0209, +0.0220]$; \code{stack\_pop} $-$ \code{super\_llm+pop} $=+0.0093\ [+0.0087, +0.0099]$. v20c so với hệ tốt nhất trước (v19c): $+0.0018\ [+0.0015, +0.0022]$, lớn hơn nhiễu seed stacker ($\pm0.0008$).}
\column{0.41\textwidth}
\footnotesize
\begin{itemize}
  \setlength{\itemsep}{3pt}
  \item Nhóm đặc trưng \code{pop} cập nhật số lượt click từ các impression diễn ra trước đó \textbf{ngay trong ngày test} (streaming).
  \item Do đó, chế độ này \textbf{không dùng để so sánh trực tiếp} với các mô hình offline hoặc các bài báo công bố trên leaderboard.
  \item Đáng chú ý, riêng \code{popularity} (không dùng nội dung bài viết) đã đạt điểm cao hơn mean pooling trên các bộ embedding (0.3787--0.3945).
\end{itemize}
\end{columns}
\par\kern0pt% chiều cao khung tính cả phần sâu của hộp cuối
\note{Chế độ online cho thấy nếu hệ thống được dùng click thời gian thực thì điểm còn tăng mạnh. Nhưng phải tách bạch: đây là thông tin mà các mô hình offline và các bài báo không có, nên không dùng con số này để so với công bố.}
\end{frame}

%% ============================================================================ Phân rã
\begin{frame}[shrink]{Phân rã mức tăng hiệu năng: $\sim$80\% từ embedding và reader, 20\% từ exposure}
\begin{columns}[T,onlytextwidth]
\column{0.53\textwidth}
\centering
\begin{tikzpicture}[font=\tiny,x=0.79cm,y=42cm]
\def\bz{0.37}
\foreach \v/\lab in {0.38/0.38,0.40/0.40,0.42/0.42,0.44/0.44}
  {\draw[black!12] (-0.15,\v-\bz)--(7.85,\v-\bz); \node[anchor=east] at (-0.15,\v-\bz) {\lab};}
\fill[black!40] (0,0) rectangle (0.75,0.3787-\bz);   \node[above] at (0.375,0.3787-\bz) {0.3787};
\fill[okc!80]   (1,0.3787-\bz) rectangle (1.75,0.3916-\bz); \node[above] at (1.375,0.3916-\bz) {+0.0130};
\fill[partc!85] (2,0.3916-\bz) rectangle (2.75,0.4102-\bz); \node[above] at (2.375,0.4102-\bz) {+0.0186};
\fill[okc!80]   (3,0.4102-\bz) rectangle (3.75,0.4260-\bz); \node[above] at (3.375,0.4260-\bz) {+0.0158};
\fill[keyc!70]  (4,0.4260-\bz) rectangle (4.75,0.4373-\bz); \node[above] at (4.375,0.4373-\bz) {+0.0114};
\fill[keyc]     (5,0) rectangle (5.75,0.4373-\bz);   \node[above] at (5.375,0.4373-\bz) {0.4373};
\fill[accentc]  (6,0.4373-\bz) rectangle (6.75,0.4516-\bz); \node[above] at (6.375,0.4516-\bz) {+0.0143};
\fill[accentc!70!black] (7,0) rectangle (7.75,0.4516-\bz); \node[above] at (7.375,0.4516-\bz) {0.4516};
\foreach \x/\lab in {0.375/mốc,1.375/+ LSE\\(H8),2.375/+ chuyên\\mục (F),3.375/+ reader\\(H10),4.375/+ kết\\hợp,5.375/\code{super}\\\code{\_zs},6.375/+ GBDT\\exposure,7.375/\code{stack}\\\code{\_meta}}
  \node[below,align=center,text width=0.86cm] at (\x,0) {\lab};
\end{tikzpicture}
\column{0.44\textwidth}
\scriptsize
\renewcommand{\arraystretch}{1.1}
\setlength{\tabcolsep}{3pt}
\begin{tabularx}{\linewidth}{@{}YR{3.65cm}@{}}
\toprule
\textbf{Bước tích hợp} & \textbf{$\Delta$ nDCG@10 [CI 95\%]}\\
\midrule
+ LSE & $+0.0130\ [+0.0123, +0.0136]$\\
+ chuyên mục & $+0.0186\ [+0.0179, +0.0192]$\\
+ reader & $+0.0158\ [+0.0150, +0.0165]$\\
+ kết hợp tuyến tính & $+0.0114\ [+0.0109, +0.0118]$\\
\midrule
tổng 4 bước nội dung & $+0.0586$ (80\%)\\
GBDT + exposure & $+0.0143\ [+0.0137, +0.0149]$\\
\bottomrule
\end{tabularx}
\end{columns}
\begin{itemize}
  \setlength{\itemsep}{1pt}
  \scriptsize
  \item Bản thân GBDT khi chỉ học trên điểm embedding \textbf{không vượt} kết hợp tuyến tính: \code{stack\_llm} $-$ \code{super\_zs} $=-0.0004\ [-0.0010, +0.0002]$.
  \item Đóng góp thực sự của GBDT đến từ nhóm đặc trưng \textbf{exposure}: \code{stack\_meta} $-$ \code{stack\_llm} $=+0.0147\ [+0.0142, +0.0153]$.
\end{itemize}
\ghichu{Mốc cơ sở: \code{qwen3:frozen}; tổng mức tăng +0.0729. Reader: \code{cat:um\_knn\_rag} (trung bình 2 seed). Số liệu từ lần chạy v16b; tỷ trọng đóng góp của từng bước phụ thuộc vào thứ tự tích hợp.}
\par\kern0pt% chiều cao khung tính cả phần sâu của hộp cuối
\note{Biểu đồ thác nước cho thấy mức tăng đến từ đâu. Khoảng 80\% mức tăng đến từ các bước về embedding và reader: LSE, chuyên mục, reader, hợp nhất. 20\% còn lại đến từ GBDT, nhưng phân tích kỹ cho thấy GBDT tự nó không hơn hợp nhất tuyến tính; phần tăng thật đến từ đặc trưng exposure.}
\end{frame}

%% ============================================================================ Split gain
\begin{frame}[shrink]{Độ quan trọng đặc trưng trong GBDT: tỷ trọng split gain trên 500 cây}
\scriptsize
Tỷ trọng độ lợi phân nhánh (split gain) theo từng nhóm đặc trưng, tổng hợp trên toàn bộ cây XGBoost:\par\vspace{2pt}
\renewcommand{\arraystretch}{1.15}
\begin{tabularx}{\textwidth}{@{}lccccY@{}}
\toprule
\textbf{Biến thể} & \textbf{\code{scores}} & \textbf{\code{ctx}} & \textbf{\code{meta}} & \textbf{\code{pop}} & \textbf{Các đặc trưng có split gain cao nhất}\\
\midrule
\code{stack\_llm} & 66\% & 34\% & -- & -- & \code{cat:um\_knn\_rag}, \code{ctx\_n\_cand}, \code{ctx\_slate\_max}, \code{bge:um\_frozen}\\
\textbf{\code{stack\_meta}} & \textbf{53\%} & \textbf{27\%} & \textbf{20\%} & -- & \code{cat:um\_knn\_rag}, \code{ctx\_n\_cand}, \code{bge:um\_frozen}, \code{ctx\_slate\_max}, \code{m\_ctr}\\
\code{stack\_pop} & 46\% & 22\% & 8\% & 24\% & \code{ctx\_n\_cand}, \code{cat:um\_knn\_rag}, \code{pop}, \code{popw\_ctr\_1h}, \code{bge:um\_frozen}\\
\bottomrule
\end{tabularx}
\par\vspace{6pt}
\centering
\begin{tikzpicture}[font=\scriptsize]
\fill[keyc] (0,0) rectangle (5.3,0.42); \node[white,font=\scriptsize\bfseries] at (2.65,0.21) {\code{scores}: 53\%};
\fill[okc!80] (5.3,0) rectangle (8.0,0.42); \node[white,font=\scriptsize\bfseries] at (6.65,0.21) {\code{ctx}: 27\%};
\fill[accentc] (8.0,0) rectangle (10.0,0.42); \node[white,font=\scriptsize\bfseries] at (9.0,0.21) {\code{meta}: 20\%};
\draw[decorate,decoration={brace,mirror,amplitude=4pt},thick,keyc] (0,-0.08)--(8.0,-0.08)
  node[midway,below=5pt,font=\scriptsize\bfseries,keyc] {80\% split gain của \code{stack\_meta} đến từ điểm mô hình nội dung và ngữ cảnh slate};
\end{tikzpicture}
\par\vspace{2pt}
\raggedright
\begin{itemize}
  \setlength{\itemsep}{1pt}
  \item Đặc trưng đơn lẻ đóng góp lớn nhất là điểm của reader \code{cat:um\_knn\_rag} (9--13\%). Với những bài đã hiển thị nhiều nhưng ít được click, đặc trưng \code{m\_ctr} giúp hạ thứ hạng bài đó.
  \item Số lượng ứng viên trong slate (\code{ctx\_n\_cand}) nằm trong nhóm đầu ở cả ba biến thể; loại bỏ đặc trưng này làm giảm điểm holdout. Cơ chế tác động cụ thể của nó cần được phân tích thêm (split gain chỉ phản ánh cấu trúc cây đã học, không mang ý nghĩa nhân quả).
\end{itemize}
\par\kern0pt% chiều cao khung tính cả phần sâu của hộp cuối
\note{Split gain cho biết GBDT dựa vào đâu. Khoảng 80\% đến từ điểm embedding và ngữ cảnh, tức là các thành phần LLM vẫn là xương sống. Kích thước slate quan trọng một cách khó giải thích; nhóm ghi nhận đây là điểm cần kiểm tra thêm, không diễn giải quá mức.}
\end{frame}

%% ============================================================================ So với công bố
\begin{frame}[shrink]{So sánh với các kết quả đã công bố trên MIND-large}
\scriptsize
\renewcommand{\arraystretch}{1.16}
\begin{tabularx}{\textwidth}{@{}lYccc@{}}
\toprule
\textbf{Mô hình} & \textbf{Nguồn / thiết lập đánh giá} & \textbf{AUC} & \textbf{MRR} & \textbf{nDCG@10}\\
\midrule
NRMS (GloVe-300) & trích dẫn từ các bài báo & 0.670--0.691 & -- & 0.418--0.434\\
PLM-NR & Wu và cs., SIGIR 2021 & 0.7064 & 0.3539 & 0.4438\\
UNBERT & Zhang và cs., IJCAI 2021 & 0.7068 & 0.3568 & 0.4478\\
Fastformer + PLM-NR & top leaderboard, số liệu thứ cấp & 0.7268 & -- & --\\
\midrule
\code{super\_zs} (chỉ dùng nội dung + lịch sử) & đề xuất, offline & 0.6863 & -- & 0.4373\\
\textbf{\code{stack\_meta}} (+ exposure) & \textbf{đề xuất, offline} & \textbf{0.7013} & \textbf{0.3503} & \textbf{0.4516}\\
\code{stack\_pop} (+ click trong ngày) & đề xuất, online & 0.7196 & 0.3671 & 0.4731\\
\midrule
\code{super\_llm} (v20c, chỉ nội dung + lịch sử) & mới nhất, offline & 0.6939 & 0.3418 & 0.4407\\
\code{stack\_meta} (v20c, + exposure) & mới nhất, offline & 0.7038 & 0.3504 & 0.4534\\
\code{stack\_pop} (v20c, + click trong ngày) & mới nhất, online & 0.7216 & 0.3676 & 0.4749\\
\bottomrule
\end{tabularx}
\begin{alertblock}{Lưu ý khi đối chiếu số liệu}
\scriptsize
\begin{itemize}
  \setlength{\itemsep}{1pt}
  \item Khác biệt tập đánh giá: nhóm đo trên \code{MINDlarge\_dev}; trong khi số liệu công bố đo trên tập test ẩn của benchmark.
  \item Hệ thống chỉ dùng nội dung và lịch sử đọc (\code{super\_zs}, \code{stack\_llm}: AUC $\approx$ 0.686; bản huấn luyện dài hơn ở v19a đạt AUC 0.692, nDCG@10 0.4420; bản mới nhất v20c: AUC 0.6939, nDCG@10 0.4407) nằm ở cận trên của NRMS nhưng vẫn thấp hơn PLM-NR khoảng 0.015--0.02 AUC.
  \item \code{stack\_meta} tiệm cận PLM-NR về AUC và cao hơn về nDCG@10, nhưng có sử dụng thêm đặc trưng exposure nên \textbf{không cùng điều kiện đầu vào}. \code{stack\_pop} dùng tín hiệu click trong ngày nên không so sánh trực tiếp.
\end{itemize}
\end{alertblock}
\par\kern0pt% chiều cao khung tính cả phần sâu của hộp cuối
\note{Khi so với công bố, nhóm nói rất cẩn thận. Nếu chỉ dùng nội dung và lịch sử, hệ của nhóm nằm trong dải của NRMS và còn kém PLM-NR khoảng 0.02 AUC. stack\_meta vượt PLM-NR về nDCG@10, nhưng nhờ thông tin exposure mà PLM-NR không dùng. Khoảng trống còn lại so với văn liệu nằm ở mô hình nội dung được tinh chỉnh end-to-end.}
\end{frame}

%% ============================================================================ Triển khai
\begin{frame}[shrink]{Khả năng triển khai: lưu cache vector + GBDT so với PLM end-to-end}
\scriptsize
\renewcommand{\arraystretch}{1.12}
\begin{tabularx}{\textwidth}{@{}P{2.2cm}YY@{}}
\toprule
\textbf{Tiêu chí} & \textbf{PLM tinh chỉnh end-to-end (PLM-NR, UNBERT)} & \textbf{Kiến trúc đề xuất (embedding cố định + GBDT)}\\
\midrule
Khi có bài báo mới & mã hóa bằng mô hình đang tinh chỉnh; cần cập nhật lại vector khi huấn luyện lại & mã hóa mỗi bài \textbf{một lần duy nhất}: $\approx$4 ms (BGE-large) đến 10 ms (Qwen3) trên T4\\
Khi phục vụ (serving) & chạy mạng BERT/RoBERTa trên GPU & tính tích vô hướng trên vector lưu sẵn + duyệt cây GBDT\\
Độ trễ phục vụ & phụ thuộc kích thước PLM và tài nguyên GPU & \textbf{chưa đo end-to-end}; ước tính vài ms cho mỗi slate trên CPU\\
Cold AUC & phụ thuộc vào mô hình đã tinh chỉnh & 0.7248, không cần cập nhật trọng số encoder\\
nDCG@10 & 0.4438--0.4478 (số liệu công bố) & 0.4516 offline / 0.4731 online (v20c: 0.4534 / 0.4749)\\
\bottomrule
\end{tabularx}
\begin{alertblock}{Ý nghĩa thực tiễn và giới hạn hiện tại}
\scriptsize Toàn bộ quy trình thực nghiệm MIND-large (14.085.557 cặp ứng viên test) chạy trọn vẹn trên 1 GPU T4 16 GB; mức RAM đỉnh ghi nhận là 21.2 GB (yêu cầu hệ thống $\ge$ 24 GB).
Chi phí độ trễ khi phục vụ mới dừng ở mức ước tính lý thuyết; cần đo benchmark end-to-end trong môi trường thực tế.
\end{alertblock}
\par\kern0pt% chiều cao khung tính cả phần sâu của hộp cuối
\note{Ưu điểm thực tế của kiến trúc là chi phí: mỗi bài chỉ mã hóa một lần, lúc phục vụ chỉ cần tích vô hướng và GBDT. Cả quy trình chạy trên một GPU T4. Nhóm chưa đo độ trễ đầu-cuối nên không khẳng định con số độ trễ; đó là việc cần làm tiếp.}
\end{frame}

%%%%%%%%%%%%%%%%%%%%%%%%%%%%%%%%%%%%%%%%%%%%%%%%%%%%%%%%%%%%%%%%%%%%%%%%%%%%%%%
\section{Giới hạn và kết luận}
%% ============================================================================ Giới hạn
\begin{frame}[shrink]{Giới hạn của nghiên cứu}
\scriptsize
\begin{enumerate}
  \setlength{\itemsep}{3pt}
  \item \textbf{Tham số trọng số chọn mô hình:} hệ số $\tau=0.80$ trên tập validation có trọng số được đặt gần với tỷ lệ click vào bài mới trên tập test (82.25\%). Trong thực tế, cần ước lượng tỷ lệ này trực tiếp từ phân phối ứng viên mà không nhìn vào nhãn test.
  \item \textbf{Phạm vi bộ dữ liệu:} chỉ đánh giá trên MIND (văn bản ngắn, chỉ có tiêu đề và tóm tắt) và dùng tập \code{*\_dev} làm tập test; số liệu của các mô hình công bố là số thứ cấp.
  \item \textbf{Quy mô mô hình ngôn ngữ:} bộ mã hóa 4B mới thử nghiệm làm đầu vào cho reader (v18a: kết quả ngang bản 0.6B); mô hình LLM sinh văn bản lớn nhất mới ở mức 1.7B, chưa thử nghiệm các mô hình 7B--8B+.
  \item \textbf{Ngân sách tinh chỉnh và kích thước mẫu:} H2 chỉ chạy trên MIND-small; User Tower LoRA (H5) mới học tối đa 3.5\% dữ liệu (1 seed); H12 mới chạy khoảng 0.5 epoch; H9 và hồ sơ của H6 chỉ đánh giá trên tập con 2.000--2.500 impression (mở rộng bài của H6 đã chạy toàn bộ ở v20a); chưa thực hiện fine-tune toàn phần trên MIND-large.
  \item \textbf{Các kết quả không hiệu quả:} phần lớn được kết luận từ bước sàng lọc trên MIND-small với một cấu hình chuẩn; riêng SID-GPT ngang ngẫu nhiên trên MIND-small do cách chấm cộng cả tầng mã mịn; sửa cách chấm trên MIND-large (v20b) tăng +0.048 nhưng vẫn kém mean pooling.
  \item \textbf{Phương sai theo seed:} khoảng tin cậy (CI) chỉ phản ánh sai số lấy mẫu trên tập test. Độ lệch giữa các seed đo được là 0.007--0.014 nDCG@10 (MIND-small) và 0.001--0.007 (reader trên MIND-large); nhiều cấu hình mới chỉ chạy 1 seed (kể cả stacker của v20c).
  \item \textbf{Đối chiếu với các công trình khác:} chưa tự huấn luyện lại PLM-NR, UNBERT trên cùng tập \code{MINDlarge\_dev}; chưa so sánh trực tiếp với MINER, GLORY hay ONCE.
  \item \textbf{Hạng mục chưa đo:} độ trễ phục vụ end-to-end; kết quả thực nghiệm của SID-News và nhóm T.
\end{enumerate}
\par\kern0pt% chiều cao khung tính cả phần sâu của hộp cuối
\note{Nhóm chủ động nêu giới hạn. Hai giới hạn quan trọng nhất: tham số 0.80 dùng thông tin tổng hợp từ test, và nhiều kết quả chỉ có một seed trong khi nhiễu seed khá lớn. Các kết quả âm cũng chủ yếu đo trên MIND-small với một cấu hình.}
\end{frame}

%% ============================================================================ Kết luận
\begin{frame}[shrink]{Kết luận: trả lời năm câu hỏi nghiên cứu}
\scriptsize
\renewcommand{\arraystretch}{1.25}
\begin{tabularx}{\textwidth}{@{}lY@{}}
\toprule
\textbf{RQ1} & Embedding cố định kết hợp phép gộp phù hợp đạt kết quả \textbf{tương đương} NRMS, NAML, Fastformer được huấn luyện lại (0.4094 so với NRMS refit 0.4078, CI chứa 0) và vượt Fastformer refit (+0.0061).\\
\textbf{RQ2} & Điểm nghẽn nằm ở \textbf{tầng tổng hợp lịch sử đọc}: phép gộp LSE tăng +0.011 đến +0.014 mà không cần tham số học; trong khi tăng kích thước encoder hầu như không cải thiện (BGE-large +0.0015; Qwen3-4B ngang bản 0.6B).\\
\textbf{RQ3} & Trong 14 nhóm kỹ thuật đã đo: \textcolor{okc}{\textbf{2 thành công rõ rệt}} (H8, H10), \textcolor{partc}{\textbf{2 hiệu quả có điều kiện}} (F, H3), \textcolor{failc}{\textbf{10 không vượt qua đối chứng mạnh}} — đặc biệt là các hướng dùng LLM để nén hoặc viết lại biểu diễn.\\
\textbf{RQ4} & Khi chỉ dùng điểm từ embedding, GBDT tương đương kết hợp tuyến tính ($-0.0004$); GBDT chỉ vượt lên nhờ bổ sung đặc trưng exposure (+0.0147). Chia CV theo khối thời gian giúp điểm validation khớp sát với test (độ lệch giảm từ $-0.0379 \to -0.0072$).\\
\textbf{RQ5} & Mở khóa 2 tầng cuối BGE-large chưa thu hẹp được khoảng cách (+0.0005, n.s.; Cold AUC $-0.0069$); User Tower LoRA vượt LSE nhưng không mang lại giá trị cộng thêm khi ghép vào hệ thống (trong phạm vi ngân sách $\approx$ 0.5 epoch và 3.5\% dữ liệu).\\
\bottomrule
\end{tabularx}
\par\kern0pt% chiều cao khung tính cả phần sâu của hộp cuối
\note{Tóm lại theo năm câu hỏi. Embedding cố định ngang baseline huấn luyện lại; điểm nghẽn là cách tổng hợp lịch sử; trong 14 họ chỉ 2 họ thành công rõ; GBDT chỉ có ích nhờ exposure; và tinh chỉnh một phần encoder chưa thu hẹp được khoảng cách với PLM tinh chỉnh end-to-end.}
\end{frame}

%% ============================================================================ Bài học + hướng tiếp
\begin{frame}[shrink]{Ba bài học rút ra và hướng phát triển tiếp theo}
\begin{columns}[T,onlytextwidth]
\column{0.49\textwidth}
\begin{block}{Ba bài học cốt lõi}
\footnotesize
\begin{enumerate}
  \setlength{\itemsep}{3pt}
  \item Trong bài toán tin tức, \textbf{cách tổng hợp lịch sử đọc} quan trọng hơn kích thước hay kiến trúc của bộ mã hóa.
  \item Các kỹ thuật dùng LLM để \textbf{nén hoặc viết lại} biểu diễn dễ làm mất chi tiết đặc trưng của bài mới, trong khi bài mới chiếm hơn 80\% lượng click.
  \item \textbf{Giao thức đánh giá quyết định tính đúng đắn của kết luận}: tập validation ``ấm'' và cách chia fold theo người dùng có thể tạo ra hiệu năng ảo cho những kỹ thuật không hiệu quả.
\end{enumerate}
\end{block}
\column{0.48\textwidth}
\begin{block}{Hướng nghiên cứu tiếp theo}
\footnotesize
\begin{itemize}
  \setlength{\itemsep}{3pt}
  \item Kiểm chứng trên bộ dữ liệu có toàn văn bài báo (EB-NeRD).
  \item Huấn luyện User Tower LLM (H5) và tinh chỉnh encoder trên toàn bộ dữ liệu MIND-large (hiện mới chạy 3.5\% và $\approx$ 0.5 epoch) với $\ge$ 3 seed.
  \item Ước lượng trọng số $\tau$ hoàn toàn không dùng thông tin từ tập test; đo độ trễ phục vụ end-to-end.
  \item Thực nghiệm đo kết quả cho SID-News và nhóm T; phân tích ablation riêng cho đặc trưng kích thước slate.
\end{itemize}
\end{block}
\end{columns}
\par\kern0pt% chiều cao khung tính cả phần sâu của hộp cuối
\note{Ba điều rút ra là thông điệp nhóm muốn hội đồng nhớ. Hướng tiếp theo đều xuất phát từ giới hạn đã nêu: dữ liệu có toàn văn, mô hình lớn hơn với nhiều seed hơn, giao thức không dùng thông tin test, và đo các hướng còn dang dở.}
\end{frame}

%% ============================================================================ Cảm ơn
\begin{frame}[c,noframenumbering]
\centering
{\Huge\bfseries\color{accentc}XIN TRÂN TRỌNG CẢM ƠN!\par}
\vspace{0.5cm}
{\large Câu hỏi và thảo luận\par}
\par\kern0pt% chiều cao khung tính cả phần sâu của hộp cuối
\note{Cảm ơn thầy cô và các bạn đã lắng nghe. Nhóm sẵn sàng trả lời câu hỏi; phần phụ lục có thuật ngữ, các câu hỏi thường gặp và danh mục các lần chạy.}
\end{frame}

%%%%%%%%%%%%%%%%%%%%%%%%%%%%%%%%%%%%%%%%%%%%%%%%%%%%%%%%%%%%%%%%%%%%%%%%%%%%%%%
\appendix
\section{Phụ lục}
%% ============================================================================ Thuật ngữ
\begin{frame}[shrink]{Phụ lục \textperiodcentered{} Thuật ngữ sử dụng trong bài}
\begin{columns}[T,onlytextwidth]
\column{0.49\textwidth}
\scriptsize
\begin{description}
  \setlength{\itemsep}{1pt}
  \item[Impression] một lượt hiển thị: gồm lịch sử đọc, danh sách bài ứng viên và nhãn click.
  \item[Bài mới / bài cũ] bài báo chưa từng / đã từng xuất hiện trong tập huấn luyện (cold / warm).
  \item[Mean pooling] biểu diễn người dùng bằng trung bình cộng vector các bài đã đọc.
  \item[LSE, top-k] gộp độ tương đồng theo hàm smooth-max (LogSumExp) hoặc lấy trung bình $k$ bài gần nhất (H8).
  \item[Reader] tầng mạng nhẹ có huấn luyện đọc lịch sử trên nền embedding cố định; ``theo ứng viên'' (candidate-aware) nghĩa là dùng chính ứng viên làm truy vấn attention (H10).
  \item[kNN RAG] làm giàu vector bài báo bằng tổ hợp tuyến tính với các bài lân cận gần nhất (H3).
  \item[Vector concat] nối (concatenate) vector Qwen3 và BGE-large thành vector 2048 chiều.
  \item[Zero-shot] chấm điểm trực tiếp không qua huấn luyện trên dữ liệu MIND.
\end{description}
\column{0.48\textwidth}
\scriptsize
\begin{description}
  \setlength{\itemsep}{1pt}
  \item[Kết hợp tuyến tính] cộng có trọng số các điểm thành phần đã chuẩn hóa z-score (\code{super\_*}).
  \item[GBDT stacker] mô hình cây quyết định tăng cường (XGBoost) học cách kết hợp điểm thành phần và đặc trưng ngữ cảnh.
  \item[Exposure] thông tin về độ phủ hiển thị của bài báo trích xuất từ log (không dùng nhãn click): tuổi bài, tần suất xuất hiện.
  \item[Online / offline] có / không sử dụng tín hiệu click phát sinh trong chính ngày test.
  \item[Time-block CV] chia cross-validation theo các khối thời gian liên tiếp thay vì chia ngẫu nhiên theo người dùng.
  \item[CI bắt cặp (paired)] khoảng tin cậy của hiệu số điểm giữa hai mô hình tính trên cùng tập impression.
  \item[Bonferroni] phép hiệu chỉnh mức ý nghĩa thống kê khi kiểm định đồng thời nhiều giả thuyết.
  \item[Phương sai seed] mức dao động kết quả giữa các lần huấn luyện cùng cấu hình nhưng khác random seed.
\end{description}
\end{columns}
\par\kern0pt% chiều cao khung tính cả phần sâu của hộp cuối
\note{Slide dự phòng: dùng khi cần giải thích nhanh một thuật ngữ.}
\end{frame}

%% ============================================================================ Câu hỏi 1
\begin{frame}[shrink]{Phụ lục \textperiodcentered{} Câu hỏi thường gặp (1/2)}
\scriptsize
\begin{description}
  \setlength{\itemsep}{4pt}
  \item[Vì sao không tinh chỉnh LLM end-to-end như PLM-NR?] Nhóm đã thử nghiệm tinh chỉnh một phần theo đúng hướng của PLM-NR (H12: mở khóa 2 tầng cuối BGE-large, có nhánh đối chứng đóng băng): chỉ tăng +0.0005 (n.s.), Cold AUC giảm $-0.0069$, không vượt qua ngưỡng dừng +0.006 đặt trước. Tinh chỉnh toàn phần đòi hỏi tài nguyên vượt quá giới hạn 1 GPU T4 và được định hướng cho nghiên cứu tiếp theo.
  \item[Việc \code{stack\_meta} vượt PLM-NR có phải là so sánh công bằng không?] Không cùng điều kiện đầu vào: \code{stack\_meta} sử dụng thêm đặc trưng exposure từ log và đánh giá trên \code{MINDlarge\_dev}; trong khi số liệu PLM-NR đo trên tập test ẩn của benchmark. Nhận định chuẩn xác là: hệ thống chỉ dùng nội dung + lịch sử (0.4373) vẫn thấp hơn PLM-NR (0.4438).
  \item[Hệ số $\tau=0.80$ có làm rò rỉ thông tin từ tập test không?] Hệ số này có tham chiếu đến thống kê tổng hợp của tập test (tỷ lệ click vào bài mới), dù không dùng nhãn của từng mẫu. Đây là giới hạn đã được nêu rõ; hướng khắc phục là ước lượng $\tau$ trực tiếp từ tỷ lệ bài mới trong danh sách ứng viên mà không cần nhãn click.
  \item[Vì sao không đưa \code{stack\_pop} vào so sánh với các bài báo công bố?] Vì mô hình này cập nhật số lượt click theo luồng ngay trong ngày test — nguồn thông tin mà các mô hình offline chuẩn không được phép truy cập.
  \item[Việc Semantic ID giảm tới $-0.1011$ có phải do lỗi cài đặt?] Trên MIND-small, SID-GPT cho kết quả ngang với vector ngẫu nhiên (0.2898 so với 0.2897). Lần chạy v20b trên MIND-large cho thấy nguyên nhân là cách chấm cộng cả 3 tầng mã (tầng mịn của bài mới làm hỏng điểm): chỉ chấm tầng chủ đề tăng $+0.0483$, mô hình có dùng lịch sử (đổi lịch sử sang người khác giảm $-0.0626$), nhưng vẫn kém mean pooling $-0.0314$. Kết luận: không phải lỗi cài đặt; mã rời rạc làm mất chi tiết của từng bài.
\end{description}
\par\kern0pt% chiều cao khung tính cả phần sâu của hộp cuối
\note{Slide dự phòng cho phần hỏi đáp.}
\end{frame}

%% ============================================================================ Câu hỏi 2
\begin{frame}[shrink]{Phụ lục \textperiodcentered{} Câu hỏi thường gặp (2/2)}
\scriptsize
\begin{description}
  \setlength{\itemsep}{4pt}
  \item[Khoảng tin cậy (CI) có bao gồm phương sai do huấn luyện không?] Không. CI bắt cặp chỉ đo sai số lấy mẫu trên tập test; phương sai giữa các seed huấn luyện được đo riêng (0.007--0.014 trên MIND-small, 0.001--0.007 cho reader trên MIND-large). Những chênh lệch nhỏ hơn biên độ này giữa hai mô hình có huấn luyện không được dùng để kết luận.
  \item[Tại sao không dùng các mô hình embedding lớn hơn (4B--8B)?] Nhóm đã thử Qwen3-Embedding-4B trên MIND-large (v18a): khi làm đầu vào cho reader, kết quả chỉ ngang bản 0.6B ($+0.0006$) và thấp hơn BGE-large ($-0.005$ đến $-0.009$); bản 8B vượt quá bộ nhớ GPU. Thực nghiệm xuyên suốt cho thấy encoder lớn hơn không đồng nghĩa với kết quả tốt hơn (BGE-large chỉ hơn BGE-small +0.0015).
  \item[Vì sao GBDT không vượt kết hợp tuyến tính khi chỉ dùng điểm embedding?] Bốn trọng số tuyến tính đã khai thác hết thông tin từ các đầu điểm thành phần ($-0.0004$, n.s.); GBDT chỉ phát huy tác dụng khi được cung cấp thêm nhóm đặc trưng phi tuyến mới như exposure (+0.0147).
  \item[LSE đã được so sánh trực tiếp với các mô hình attention học được chưa?] Đã so sánh. Trên cùng bảng vector cố định: đã đối chiếu với reader kiến trúc NRMS, Fastformer, đa truy vấn ($K=4$) và theo ứng viên (xem slide ``H8 so với các tầng attention có huấn luyện''). Đối với các baseline đầy đủ (NRMS, NAML, Fastformer học từ đầu): xem slide RQ1.
  \item[Việc kích thước slate đứng đầu về split gain có hợp lý không?] Đây là điểm cần phân tích thêm. Loại bỏ cột này làm giảm điểm holdout; có khả năng cây quyết định dùng độ dài danh sách ứng viên làm biến điều kiện để điều chỉnh phân phối điểm. Cần một nghiên cứu ablation riêng cho hiện tượng này.
  \item[Đóng góp chính của đề tài là gì?] (1) Đánh giá thực nghiệm 12 nhóm kỹ thuật LLM 2024--2026 trên bài toán tin tức với baseline được huấn luyện lại chuẩn xác, có CI bắt cặp và kiểm soát phương sai seed; (2) Xác định điểm nghẽn nằm ở tầng tổng hợp lịch sử và đề xuất hai giải pháp hiệu quả trên embedding cố định (LSE, reader theo ứng viên); (3) Thiết lập giao thức đánh giá chống lệch phân phối (validation có trọng số bài mới + Time-block CV); (4) Xây dựng hệ thống 3 tầng đạt hiệu năng cao và chạy trọn vẹn trên 1 GPU T4 phổ thông.
\end{description}
\par\kern0pt% chiều cao khung tính cả phần sâu của hộp cuối
\note{Slide dự phòng cho phần hỏi đáp.}
\end{frame}

%% ============================================================================ Các lần chạy
\begin{frame}[shrink]{Phụ lục \textperiodcentered{} Danh mục các lần chạy và nguồn số liệu}
\scriptsize
\renewcommand{\arraystretch}{1.18}
\begin{tabularx}{\textwidth}{@{}llYl@{}}
\toprule
\textbf{Lần chạy} & \textbf{Dữ liệu} & \textbf{Nội dung thực nghiệm chính} & \textbf{Thời gian}\\
\midrule
V10 & MIND-small & bộ đo ban đầu: NRMS, NAML, Fastformer, BGE zero-shot, popularity & --\\
v7 & MIND-small & khảo sát $\sim$16 nhóm kỹ thuật, $\sim$150 biến thể; huấn luyện lại baseline chuẩn; ablation từng bước NRMS & 7.52 giờ\\
v13 & MIND-small & khảo sát trường văn bản (F), lịch sử 100/200 click, đa sở thích (H11), Fastformer 2 seed + refit & 8.22 giờ\\
v14 & MIND-large & chạy một encoder (Qwen3) với đầy đủ kỹ thuật; kết hợp tuyến tính; User Tower LoRA (H5) & 2.90 giờ\\
v15 & MIND-large & tích hợp GBDT stacker (XGBoost GPU), đặc trưng exposure, bổ sung encoder thứ hai BGE-large & 2.73 giờ\\
v16b & MIND-large & áp dụng Time-block CV, huấn luyện reader trên BGE-large và vector concat, đo RAM đỉnh (21.2 GB) & 4.49 giờ\\
v17 & MIND-large & mở khóa 2 tầng cuối BGE-large dưới reader (H12), chạy nhánh đối chứng bắt cặp và kiểm tra ngưỡng dừng & --\\
v18a & MIND-large & tái lập kết quả v16b (khớp đến chữ số thập phân thứ 4); thử nghiệm Qwen3-Embedding-4B làm encoder & 8.15 giờ\\
v18b & MIND-large & huấn luyện reader trên toàn bộ 2.71 triệu mẫu trong 6 epoch; thử nghiệm biến thể reader \code{sca} & 5.13 giờ\\
v19a & MIND-large & đưa các reader huấn luyện dài hơi vào bộ kết hợp tuyến tính + GBDT; đo H6, H9 trên tập con & 9.54 giờ\\
v19b & MIND-large & huấn luyện H5 trong 200 phút, chấm trên toàn bộ tập test; đo H1 trên 60.000 impression & 10.09 giờ\\
v19c & MIND-large & ghép reader huấn luyện dài hơi (v19a) và H5 (v19b) vào bộ kết hợp + GBDT & --\\
v20a & MIND-large & H6: Qwen3-1.7B viết mở rộng cho toàn bộ 104.151 bài; reader trên vector có đoạn LLM viết & 5.35 giờ\\
v20b & MIND-large & H7: SID chấm theo tầng; H1 + MLP học được (\code{histq\_mlp}), toàn bộ tập test & 3.2 giờ\\
v20c & MIND-large & ghép thành phần của v20a, v20b vào hệ tốt nhất (v19c); mô hình mới nhất & 3.72 giờ\\
\bottomrule
\end{tabularx}
\ghichu{Mỗi lần chạy là một phiên notebook ``Run All'' hoàn chỉnh trên Kaggle (GPU Tesla T4 16 GB; v14--v17 dùng 1 GPU, v18--v19 dùng 2 GPU). Các bảng kết quả MIND-large trong bài lấy từ v16b, ngoại trừ các dòng đánh dấu (*) lấy từ v17 và các số liệu ghi chú rõ từ v18, v19. Nhánh đối chứng của v18a, v18b, v19a, v19b khớp hoàn toàn với v16b đến 4 chữ số thập phân; đối chứng của v20c khớp v19c.}
\par\kern0pt% chiều cao khung tính cả phần sâu của hộp cuối
\note{Slide dự phòng: cho biết mỗi con số đến từ lần chạy nào.}
\end{frame}

%% ============================================================================ Tài liệu
\begin{frame}[shrink]{Phụ lục \textperiodcentered{} Tài liệu tham khảo chính}
\tiny
\begin{columns}[T,onlytextwidth]
\column{0.49\textwidth}
\begin{itemize}
  \setlength{\itemsep}{1pt}
  \item F. Wu và cs. MIND: A Large-scale Dataset for News Recommendation. ACL 2020.
  \item C. Wu và cs. Neural News Recommendation with Multi-Head Self-Attention (NRMS). EMNLP-IJCNLP 2019.
  \item C. Wu và cs. Neural News Recommendation with Attentive Multi-View Learning (NAML). IJCAI 2019.
  \item C. Wu và cs. Fastformer: Additive Attention Can Be All You Need. arXiv 2021.
  \item C. Wu và cs. Empowering News Recommendation with Pre-trained Language Models (PLM-NR). SIGIR 2021.
  \item Q. Zhang và cs. UNBERT: User-News Matching BERT for News Recommendation. IJCAI 2021.
  \item T. Qi và cs. News Recommendation with Candidate-aware User Modeling (CAUM). SIGIR 2022.
  \item J. Li và cs. MINER: Multi-Interest Matching Network for News Recommendation. Findings of ACL 2022.
  \item Y. Cen và cs. Controllable Multi-Interest Framework for Recommendation (ComiRec). KDD 2020.
  \item O. Khattab, M. Zaharia. ColBERT. SIGIR 2020.
  \item A. Iana và cs. Simplifying Content-Based Neural News Recommendation. SIGIR 2023.
\end{itemize}
\column{0.49\textwidth}
\begin{itemize}
  \setlength{\itemsep}{1pt}
  \item P. BehnamGhader và cs. LLM2Vec. COLM 2024. \quad Q. Liu và cs. LLMEmb. AAAI 2025.
  \item Y. He và cs. LLM2Rec. KDD 2025. \quad Y. Zhang và cs. Qwen3 Embedding. arXiv 2025.
  \item C. Zhang và cs. EmbSum. RecSys 2024.
  \item Y. Xi và cs. Towards Open-World Recommendation with Knowledge Augmentation from LLMs (KAR). RecSys 2024.
  \item Q. Liu và cs. ONCE. WSDM 2024. \quad Y. Yada, H. Yamana. Category Description by an LLM. arXiv 2024.
  \item D. Edge và cs. From Local to Global: A Graph RAG Approach. arXiv 2024.
  \item S. Wang và cs. Knowledge Graph RAG for LLM-based Recommendation. ACL 2025.
  \item S. Rajput và cs. Recommender Systems with Generative Retrieval (TIGER). NeurIPS 2023.
  \item W. Wang và cs. Learnable Item Tokenization for Generative Recommendation (LETTER). CIKM 2024.
  \item Y. Hou và cs. Generating Long Semantic IDs in Parallel (RPG). KDD 2025.
  \item Y. Hou và cs. LLMs are Zero-Shot Rankers for Recommender Systems. ECIR 2024. \quad J. Gao và cs. LLM4Rerank. WWW 2025.
  \item Y. Wang và cs. Pre-train, Align, and Disentangle (PAD). SIGIR 2025.
  \item J. Kruse và cs. EB-NeRD: A Large-Scale Dataset for News Recommendation. RecSys Challenge 2024.
\end{itemize}
\end{columns}
\par\kern0pt% chiều cao khung tính cả phần sâu của hộp cuối
\note{Slide dự phòng: danh sách tài liệu đã được đối chiếu tên, tác giả và hội nghị.}
\end{frame}

\end{document}
